% 第4章 Skill体系：赋予AI金融专业能力
% 智慧银行实验教程——AI驱动的金融科技实践

\chapter{第4章\quad Skill体系：赋予AI金融专业能力}
\chapteroverview{
本章学习目标：
\begin{enumerate}
  \item 理解Skill的设计原理与知识工程视角
  \item 熟悉本机 167 个 Skill 的 20 个分类体系，重点掌握 11 类高关联度 Skill 的能力边界
  \item 掌握金融类Skill的调用方法（含 8 个从 \textbullet 到 \textbullet\textbullet\textbullet\textbullet 的实战任务）
  \item 掌握学术与商业图形展示 Skill 体系（\texttt{econ-visualization} / \texttt{paper-figure} / \texttt{paper-illustration} / \texttt{mermaid-diagram}）
  \item 掌握演示文稿 Skill 体系（学术线 \texttt{beamer-presentation} / \texttt{paper-slides} / \texttt{paper-poster}；商业线 \texttt{pptx} / \texttt{qoderwork-ppt} / \texttt{ppt-master}）
  \item 掌握元 Skill 体系（\texttt{create-skill} / \texttt{find-skills} / \texttt{skill-auditor} / \texttt{write-skill-readme} / \texttt{install-skill-dependency} / \texttt{luban}）
  \item 独立编写符合规范的金融Skill
  \item 从银行合规视角进行Skill安全审计
\end{enumerate}
}

% ===========================================================
\section{Skill库全景图：从167个Skill看AI能力生态}
% ===========================================================

本章Skill体系不是一个孤立的技术点，而是构筑在一个涵盖\textbf{167个Skill}、分\textbf{20个分类}的完整Skill生态之上。本节先绘制一张全景图，让你在深入具体Skill之前，先看清AI能力生态的宏观面貌。

\begin{figure}[H]
\centering
\begin{tikzpicture}[
  box/.style={rectangle, rounded corners=4pt, draw=blue!70, very thick, fill=blue!8, minimum width=2.6cm, minimum height=1.6cm, align=center, font=\footnotesize, inner sep=4pt},
  boxhot/.style={rectangle, rounded corners=4pt, draw=orange!80, very thick, fill=orange!15, minimum width=2.6cm, minimum height=1.6cm, align=center, font=\footnotesize, inner sep=4pt},
  arrow/.style={-{Stealth[length=3mm]}, very thick, blue!60},
]
% 4 象限：业务阶段 × 使用频率
\node[box]    (q1) at (-3.5, 1.8)  {\textbf{高频主流程}\\\textcolor{blue!70!black}{01-ideation}\\\textcolor{blue!70!black}{02-literature}\\\textcolor{blue!70!black}{03-theory}};
\node[boxhot] (q2) at ( 3.5, 1.8)  {\textbf{高频主流程}\\\textcolor{orange!80!black}{07-writing}\\\textcolor{orange!80!black}{10-communication}\\\textcolor{orange!80!black}{12-finance}};
\node[box]    (q3) at (-3.5,-1.8)  {\textbf{低频辅助}\\\textcolor{blue!70!black}{16-office}\\\textcolor{blue!70!black}{13-consulting}\\\textcolor{blue!70!black}{20-messaging}};
\node[boxhot] (q4) at ( 3.5,-1.8)  {\textbf{低频辅助}\\\textcolor{orange!80!black}{19-devtools}\\\textcolor{orange!80!black}{15-browser}\\\textcolor{orange!80!black}{17-ima}};

% 中心标签
\node[draw=red!70, very thick, fill=yellow!20, rounded corners=3pt, font=\small\bfseries, align=center, minimum width=2.2cm, minimum height=1.4cm] at (0,0) {167\\个\\Skill\\\footnotesize 20 分类};

% 分类轴
\draw[arrow, red!70] (-5.5, 0) -- (5.5, 0) node[right, font=\footnotesize\bfseries, black] {金融专业性};
\node[font=\footnotesize\bfseries, red!70] at (-5.8, 0.3) {通用};
\node[font=\footnotesize\bfseries, red!70] at (5.8, 0.3) {金融专业};

\draw[arrow, green!60!black] (0, -3.0) -- (0, 3.0) node[above, font=\footnotesize\bfseries, black] {使用频率};
\node[font=\footnotesize\bfseries, green!60!black] at (0.3, -3.2) {低};
\node[font=\footnotesize\bfseries, green!60!black] at (0.3, 3.2) {高};
\end{tikzpicture}
\caption{Skill 生态四象限全景图（自制）}
\label{fig:skill-quadrants}
\end{figure}

\subsection{本机Skill库总览}

本课程使用的Skill库位于\texttt{qoder-skills-library/}目录（CNB仓库：\url{https://cnb.cool/xiaosicau/qoder-skills-library.git}），共\textbf{167个Skill / 20个分类}，外加\texttt{.agents/skills/}项目本地的9个核心Skill。

\begin{table}[H]
\centering
\begin{tabular}{lcp{9cm}}
\toprule
\textbf{分类} & \textbf{数量} & \textbf{与银行金融业务关联度} \\
\midrule
01-ideation 选题 & 6 & 高（选题辅助） \\
02-literature 文献综述 & 10 & 高（实证文献支撑） \\
03-theory 理论建模 & 5 & 中（计量模型推导） \\
\rowcolor{yellow!20}04-data 数据获取与清洗 & \textbf{3} & \textbf{极高} \\
\rowcolor{yellow!20}05-analysis 统计分析与仿真 & \textbf{11} & \textbf{极高} \\
06-experiment 实验流程 & 5 & 中 \\
\rowcolor{yellow!20}07-writing 论文写作 & \textbf{7} & \textbf{高} \\
08-funding NSFC 申报 & 6 & 中（项目申报） \\
09-review 审稿 & 11 & 高（成果改进） \\
\rowcolor{yellow!20}10-communication 演示展示 & \textbf{7} & \textbf{高} \\
11-perspective 思维顾问 & 2 & 高（中国金融学体系、Munger/Taleb等思维） \\
\rowcolor{yellow!20}12-finance 金融与商业 & \textbf{11} & \textbf{极高} \\
13-consulting 咨询框架 & 3 & 高 \\
14-cnb CNB 平台 & 13 & 低 \\
15-browser 浏览器自动化 & 9 & 中 \\
\rowcolor{yellow!20}16-office 文档办公 & \textbf{8} & \textbf{高} \\
17-ima IMA 笔记 & 3 & 低 \\
18-opencli 站点CLI & 6 & 低 \\
\rowcolor{yellow!20}19-devtools 元工具 & \textbf{20} & \textbf{高（含skill-auditor/create-skill/find-skills等元Skill）} \\
20-messaging 通讯 & 4 & 低 \\
\bottomrule
\end{tabular}
\end{table}

其中，\textbf{11类高关联度Skill}（标黄）是本章重点。

\subsection{学术研究全流程与Skill映射}

一个完整的学术研究项目可以划分为：选题→文献→理论→数据→分析→实验→写作→审稿→展示九大环节。每一个环节都对应多个Skill可以辅助：

\begin{table}[H]
\centering
\begin{tabular}{lp{6cm}p{6cm}}
\toprule
\textbf{研究环节} & \textbf{典型任务} & \textbf{可调用的Skill} \\
\midrule
选题 & 从热点问题中寻找创新点 & \texttt{idea-generation} / \texttt{research-refine} / \texttt{ideation-spark} \\
文献综述 & 系统检索+综述+管理引用 & \texttt{find-literature} / \texttt{systematic-literature-review} / \texttt{lit-review-assistant} / \texttt{pdf-extraction} \\
理论建模 & 构建数理模型与推导 & \texttt{theory-modeling} / \texttt{model-derivation} \\
数据获取 & 从API/数据库拉取金融数据 & \texttt{tushare-finance} / \texttt{china-stock-analysis} / \texttt{yahoo-finance} / \texttt{worldbank-data360} \\
统计分析 & 描述统计+回归+因果推断 & \texttt{stata} / \texttt{r-econometrics} / \texttt{python-panel-data} / \texttt{data-analysis} \\
实验流程 & A/B测试与实验跟踪 & \texttt{ab-test-analysis} / \texttt{experiment-tracker} \\
论文写作 & LaTeX写作+表格+排版 & \texttt{paper-write} / \texttt{latex-tables} / \texttt{academic-paper-writer} / \texttt{humanizer-zh} \\
审稿与润色 & 多轮迭代提升论文质量 & \texttt{paper-review} / \texttt{peer-review} / \texttt{research-review} / \texttt{auto-review-loop} \\
演示展示 & 会议汇报与海报 & \texttt{beamer-presentation} / \texttt{paper-slides} / \texttt{paper-poster} / \texttt{ppt-master} \\
\bottomrule
\end{tabular}
\end{table}

\subsection{银行业务工作流与Skill映射}

从业务侧看，银行业的工作流可以分为：数据获取→分析→报告→决策→运营五大环节。每个环节同样有Skill支撑：

\begin{table}[H]
\centering
\begin{tabular}{lp{6cm}p{6cm}}
\toprule
\textbf{业务环节} & \textbf{典型任务} & \textbf{可调用的Skill} \\
\midrule
数据获取 & 从Wind/CSMAR/行内系统拉数 & \texttt{tushare-finance} / \texttt{api-data-fetcher} / \texttt{china-stock-mcp} \\
数据分析 & 信用评估/反欺诈/反洗钱/压力测试 & \texttt{finance-expert} / \texttt{r-econometrics} / \texttt{stata-regression} \\
报告生成 & Word/Excel/PPT 文档 & \texttt{docx} / \texttt{xlsx} / \texttt{pptx} / \texttt{pdf} \\
决策辅助 & 波特五力/SWOT/估值 & \texttt{consulting-frameworks} / \texttt{financial-modeling} / \texttt{stock-analysis} \\
运营管理 & 舆情监控+风险预警+客户分群 & \texttt{finance-news} / \texttt{bank-risk-alert} / \texttt{questionnaire-design} \\
\bottomrule
\end{tabular}
\end{table}

\subsection{Skill发现与生态接入}

如何在众多Skill中找到合适的？三种途径：

\begin{enumerate}[leftmargin=2em]
  \item \textbf{\texttt{find-skills} 元Skill}：从\texttt{skills.sh}市场、GitHub anthropics/skills 仓库、本地 qoder-skills-library、npm registry中发现Skill。
  \item \textbf{本机浏览}：直接查阅\texttt{qoder-skills-library/<分类>/<skill-name>/SKILL.md}（每个Skill一个独立目录）。
  \item \textbf{自然语言提示}：在AI IDE中描述需求，AI自动匹配最合适的Skill。
\end{enumerate}

\begin{lstlisting}[style=shell]
# 技能库根目录一览
ls qoder-skills-library/
# 查找所有包含 "finance" 的 Skill
grep -rl "finance" qoder-skills-library/*/SKILL.md
\end{lstlisting}

\begin{notebox}[Skill发现的三条经验]
\begin{itemize}[leftmargin=1em]
  \item \textbf{不确定时先试用\texttt{find-skills}元Skill}：它会同时查询多个生态市场，给你最全的候选。
  \item \textbf{优先选择本地Skill}：本机已安装的Skill不需要额外下载，响应更快。
  \item \textbf{看SKILL.md的tags字段}：tags是Skill的能力标签，是快速判断Skill适用场景的捷径。
\end{itemize}
\end{notebox}

% ===========================================================
\section{Skill设计原理}
% ===========================================================

\subsection{知识工程视角：从专家系统到AI Skill的演进}

人工智能在专业领域的应用经历了从"硬编码知识"到"灵活编排知识"的深刻变革。理解这一演进脉络，有助于我们把握Skill体系的设计哲学。

\subsubsection{专家系统时代（1980s--2000s）}

20世纪80年代，知识工程的核心范式是\textbf{专家系统}（Expert System）。其基本思路是将领域专家的知识以"如果--那么"（IF-THEN）规则的形式编码进计算机系统。典型代表如MYCIN（医学诊断）、XCON（计算机配置）等。

专家系统的核心架构包含三个组件：
\begin{enumerate}[leftmargin=2em]
  \item \textbf{知识库}（Knowledge Base）：存储领域规则，如"如果资产负债率$>$70\%，则风险等级为高"
  \item \textbf{推理引擎}（Inference Engine）：基于规则进行前向或后向链式推理
  \item \textbf{用户界面}：供非专业用户与系统交互
\end{enumerate}

专家系统的致命缺陷在于\textbf{知识获取瓶颈}——将专家的隐性知识显式编码为规则，既耗时又容易遗漏。一个银行信贷审批专家系统可能需要数千条规则，而任何一条遗漏都可能导致错误决策。

\subsubsection{机器学习时代（2010s）}

深度学习的崛起提供了一种全新路径：不再人工编写规则，而是让模型从数据中自动学习。然而，纯数据驱动的方法在金融领域面临独特挑战：
\begin{itemize}[leftmargin=2em]
  \item 金融数据稀缺且信噪比低，黑天鹅事件无法从历史数据中学习
  \item 监管要求模型具有\textbf{可解释性}，"黑箱"决策难以通过合规审查
  \item 业务流程中蕴含大量\textbf{流程性知识}（如"先查征信再审批"），无法从数据中自动习得
\end{itemize}

\subsubsection{AI Skill时代（2024--）}

Skill体系代表了知识工程的最新范式：\textbf{以结构化Markdown文档承载领域知识，通过自然语言指导大模型执行专业任务}。它既非硬编码的规则库，也非纯数据驱动的黑箱，而是将人类的\textbf{方法论知识}（做什么、怎么做、注意什么）以AI可理解的方式注入工作流。

\begin{theorybox}[知识工程三大范式对比]
\begin{tabular}{lp{3.5cm}p{3.5cm}p{3.5cm}}
\toprule
\textbf{维度} & \textbf{专家系统} & \textbf{机器学习} & \textbf{AI Skill} \\
\midrule
知识载体 & IF-THEN规则 & 训练数据权重 & Markdown文档 \\
知识来源 & 人工编码 & 数据学习 & 人机协作 \\
可解释性 & 强 & 弱 & 强 \\
适应性 & 差（规则固定） & 中（需重新训练） & 好（修改文档即可） \\
金融适用性 & 流程固定但难维护 & 模式识别强但不透明 & 兼顾流程与灵活性 \\
\bottomrule
\end{tabular}
\end{theorybox}

\subsection{SKILL.md规范详解}

Skill的物理形态是一个名为\texttt{SKILL.md}的Markdown文件，遵循统一的规范结构。这一规范是AI IDE（如Qoder、TRAE）识别和调用Skill的基础。

\begin{table}[H]
\centering
\begin{tabular}{lp{9cm}}
\toprule
\textbf{组成部分} & \textbf{说明} \\
\midrule
\texttt{name} / \texttt{description} & 名称与一句话描述，AI据此判断何时调用。description应简洁明确，控制在100字符以内 \\
\texttt{When to Use} & 明确的触发场景（触发词 + 使用条件），定义Skill的适用边界 \\
\texttt{Instructions} & 分步骤的操作流程（Step 1, 2, 3...），是Skill的核心执行逻辑 \\
\texttt{Example Prompts} & 示例提示词，展示典型用法，帮助AI理解预期输入格式 \\
\texttt{Requirements} & 运行依赖（Python包、CLI工具等），确保执行环境就绪 \\
\texttt{Best Practices} & 最佳实践与注意事项，防止常见错误 \\
\bottomrule
\end{tabular}
\end{table}

上述六个部分并非随意组合，而是遵循\textbf{"发现--理解--执行--保障"}的逻辑链条：

\begin{verbatim}
┌─────────────────────────────────────────────────────────────┐
│  name/description → AI判断"这是不是我该做的事"            │
│  When to Use      → AI判断"现在是不是该做这件事"          │
│  Instructions     → AI知道"具体怎么做"                    │
│  Example Prompts  → AI理解"用户会怎么叫我做"              │
│  Requirements     → AI检查"我有没有能力做"                │
│  Best Practices   → AI注意"做的时候别犯这些错"            │
└─────────────────────────────────────────────────────────────┘
\end{verbatim}

\begin{notebox}[SKILL.md的YAML前缀]
除了正文内容外，SKILL.md通常还包含一个YAML格式的前置元数据（frontmatter），用于结构化描述Skill的基本属性：
\begin{lstlisting}[style=json]
---
name: finance-expert
description: Financial analysis expert for banking
workflow_stage: analysis
compatibility:
  - claude-code
  - cursor
  - qoder
version: 1.0.0
tags:
  - finance
  - banking
  - risk-analysis
---
\end{lstlisting}
这些元数据帮助Skill市场进行索引和分类，也有助于AI IDE进行兼容性检查。
\end{notebox}

\subsection{MCP vs Skill的本质区别}

在前一章中我们学习了MCP协议，本章又引入了Skill体系。初学者常将二者混淆，但它们解决的是截然不同的问题。

\begin{theorybox}[黄金口诀：MCP给能力，Skill给方法]
\begin{verbatim}
┌─────────────────────────────────────────────────────────────┐
│  MCP = AI 的「手脚」（连接外部工具）                        │
│  Skill = AI 的「专业脑」（领域知识 + 工作流程）            │
│  两者配合 = AI 既知道怎么做，又有工具去做                  │
└─────────────────────────────────────────────────────────────┘
\end{verbatim}

更精确地说：
\begin{itemize}[leftmargin=1em]
  \item \textbf{MCP}解决的是\textbf{能力}问题——AI能否读写文件、访问数据库、调用API？MCP为AI提供了与外部世界交互的通道，但不告诉AI应该做什么。
  \item \textbf{Skill}解决的是\textbf{方法}问题——面对一个信用风险评估任务，应该分几步走？每步用什么指标？输出什么格式？Skill为AI提供了专业的方法论框架。
\end{itemize}

以银行信用风险评估为例：
\begin{itemize}[leftmargin=1em]
  \item 没有MCP：AI无法读取客户财务数据，无法生成Word报告
  \item 没有Skill：AI可以读取数据，但不知道该算哪些指标、按什么标准评级
  \item MCP + Skill：AI既能获取数据、生成报告，又知道该按五C模型进行专业评估
\end{itemize}
\end{theorybox}

\subsection{Skill的触发机制：AI如何决定何时调用}

Skill的触发是一个\textbf{语义匹配}过程。当用户在AI IDE中输入提示词时，AI会执行以下判断流程：

\begin{enumerate}[leftmargin=2em]
  \item \textbf{意图解析}：理解用户想要完成什么任务
  \item \textbf{Skill检索}：在已安装的Skill列表中搜索，比对用户意图与每个Skill的\texttt{name}/\texttt{description}/\texttt{When to Use}
  \item \textbf{相关性评分}：对匹配到的Skill进行相关性打分，选择得分最高的
  \item \textbf{参数提取}：从用户提示词中提取Skill所需的输入参数
  \item \textbf{执行与反馈}：按照\texttt{Instructions}逐步执行，将结果反馈给用户
\end{enumerate}

这一过程的核心是\textbf{触发词}（Trigger Words）的设计。好的触发词应当让AI在正确场景下精确匹配，同时避免在不相关场景下误触发。例如：
\begin{itemize}[leftmargin=2em]
  \item \texttt{finance-expert}的触发词："信用风险评估""DCF估值""DuPont分析"
  \item \texttt{bank-risk-alert}的触发词："风险预警""贷后预警""逾期提醒"
  \item 如果触发词设计为"分析"——太宽泛，任何场景都可能误触发
  \item 如果触发词设计为"使用Altman Z-Score评估恒达机械"——太具体，限制了Skill的适用范围
\end{itemize}

\subsection{Skill的生命周期}

一个Skill从创建到发挥作用，经历完整的生命周期：

\begin{verbatim}
┌─────────────────────────────────────────────────────────────┐
│                                                             │
│  安装 → 发现 → 匹配 → 执行 → 反馈                         │
│                                                             │
│  [安装]  npx skills add / 本地创建SKILL.md                  │
│  [发现]  AI扫描已安装Skill列表，建立索引                    │
│  [匹配]  用户提示词 → 语义匹配 → 选择最佳Skill              │
│  [执行]  按Instructions逐步执行，调用MCP工具               │
│  [反馈]  输出结果 + 用户评价 → 迭代优化Skill                │
│                                                             │
└─────────────────────────────────────────────────────────────┘
\end{verbatim}

其中\textbf{反馈}环节往往被忽视，但它是Skill持续优化的关键。一个好的Skill应当根据实际使用中的问题不断迭代——比如发现某个触发词容易误触发，就在\texttt{When to Use}中增加排除条件；发现某个步骤容易出错，就在\texttt{Best Practices}中增加提醒。

% ===========================================================
\section{运行环境准备详解：Skill调用前的环境检查}

% ===========================================================

\textbf{【本章特别要求】} 调用Skill不是点下按钮那么简单。许多Skill（如\texttt{r-econometrics}、\texttt{econ-visualization}、\texttt{qoderwork-ppt}）背后依赖完整的语言解释器、外部包、Node模块或API Key。学生在实验时最常见的问题是"命令找不到"、"包未安装"、"路径错误"。本节集中讲解调用Skill所需的运行环境，避免你在实验时遇到这些坑。

\subsection{R语言环境（r-econometrics / econ-visualization的前置依赖）}

\textbf{R解释器安装}（macOS推荐Homebrew）：

\begin{lstlisting}[style=shell]
# 方式1：Homebrew安装最新稳定版
brew install --cask r

# 方式2：下载官方安装包（适合离线/教学机房）
# https://cran.r-project.org/bin/macosx/  下载 R-4.x.x-arm64.pkg
\end{lstlisting}

\textbf{RStudio Desktop}（IDE，推荐）：\url{https://posit.co/download/rstudio-desktop/}

\textbf{CRAN镜像配置}（加速国内下载）：在\texttt{~/.Rprofile}中写入：

\begin{lstlisting}[style=r]
options(repos = c(CRAN = "https://mirrors.tuna.tsinghua.edu.cn/CRAN/"))
\end{lstlisting}

\textbf{必备R包一键安装}（\texttt{r-econometrics} / \texttt{econ-visualization}共需）：

\begin{lstlisting}[style=r]
install.packages(c(
  "tidyverse",   # 数据清洗+ggplot2基础
  "fixest",      # 高维固定效应+IV+DiD（替代reghdfe/ivreg2）
  "modelsummary",# 发表级回归表
  "rdrobust",    # 断点回归
  "did",         # Callaway-Sant'Anna现代DID
  "Synth",       # 合成控制法
  "MatchIt",     # 倾向得分匹配
  "sandwich", "lmtest",  # 稳健标准误
  "AER",         # IV估计（ivreg）
  "stargazer"    # 回归表替代方案
))
\end{lstlisting}

\textbf{R路径验证}：

\begin{lstlisting}[style=shell]
which R
R --version          # 应输出 R version 4.x.x
Rscript -e 'cat(packageVersion("fixest"))'   # 应输出可读版本号
\end{lstlisting}

\subsection{Stata环境（stata / stata-regression的前置依赖）}

\textbf{macOS路径定位}（学生常踩坑）：

\begin{lstlisting}[style=shell]
# Stata 18 / StataNow标准路径
/Applications/Stata/StataMP.app/Contents/MacOS/stata-mp
/Applications/StataNow/StataMP.app/Contents/MacOS/stata-mp
# 自动定位
mdfind -name "stata-mp" | grep MacOS
\end{lstlisting}

\textbf{PATH配置}（加入\texttt{~/.zshrc}）：

\begin{lstlisting}[style=shell]
export PATH="/Applications/Stata/StataMP.app/Contents/MacOS:$PATH"
source ~/.zshrc
\end{lstlisting}

\textbf{批处理模式调用}（AI Skill调用Stata的标准方式）：

\begin{lstlisting}[style=shell]
stata-mp -b do analysis.do     # 输出到 analysis.log
stata-mp -b do /tmp/run.do && echo SUCCESS || echo FAILED
\end{lstlisting}

\textbf{必备外部命令}（\texttt{ssc install}一键安装）：

\begin{lstlisting}[style=stata]
ssc install reghdfe, replace
ssc install estout, replace
ssc install outreg2, replace
ssc install did_multiplegt, replace
ssc install eventstudyinteract, replace
ssc install rdrobust, replace
ssc install psmatch2, replace
ssc install synth, replace
ssc install ivreg2, replace
ssc install xtabond2, replace
ssc install winsor2, replace
\end{lstlisting}

\subsection{Python环境（python-panel-data / data-analysis / paper-figure的前置依赖）}

\textbf{建议版本}：Python 3.10+（与\texttt{reghdfe}等Stata包无冲突）。

\textbf{Conda环境}（推荐，隔离R/Stata/Python）：

\begin{lstlisting}[style=shell]
conda create -n banking python=3.11
conda activate banking
pip install pandas numpy statsmodels scikit-learn pmdarima prophet torch \
            matplotlib seaborn linearmodels PyFixest   # PyFixest是fixest的Python版
\end{lstlisting}

\textbf{Jupyter Lab}（推荐IDE）：\texttt{pip install jupyterlab}

\subsection{PPT演示文稿相关环境（pptx / qoderwork-ppt / ppt-master的前置依赖）}

\textbf{Node.js + npm}：\url{https://nodejs.org/}（LTS版）。

\textbf{Puppeteer + LibreOffice}（\texttt{qoderwork-ppt}一键管线需要）：

\begin{lstlisting}[style=shell]
brew install --cask libreoffice   # macOS
npm install -g pptxgenjs puppeteer jsdom dom-to-pptx lucide-static
\end{lstlisting}

\textbf{poppler-utils}（\texttt{pptx} Skill的PDF转图片依赖）：

\begin{lstlisting}[style=shell]
brew install poppler   # 提供 pdftoppm 命令
\end{lstlisting}

\textbf{LaTeX Beamer}（\texttt{beamer-presentation} / \texttt{paper-slides}需要）：macTeX已包含；Linux \texttt{sudo apt install texlive-latex-extra texlive-fonts-recommended}。

\begin{warningbox}[Skill使用前必走的3步环境检查]
每个Skill被学生调用前，请先运行3步检查，避免命令找不到的尴尬：

\begin{lstlisting}[style=shell]
# === Step 1：路径/可执行文件检查 ===
which <skill-binary-or-interpreter>
# 例：which R / which stata-mp / which python3 / which xelatex

# === Step 2：依赖包/库检查 ===
<interpreter> -e "<import-test-command>"
# 例：Rscript -e 'library(fixest)' / python3 -c 'import tushare'

# === Step 3：Skill可访问性检查 ===
ls <path-to-skill>/SKILL.md
# 例：ls qoder-skills-library/05-analysis/stata/SKILL.md
\end{lstlisting}

\begin{itemize}[leftmargin=1em]
  \item 三步全过 → 可使用该Skill
  \item Step 1不过 → 安装语言解释器（brew install / 官网下载）
  \item Step 2不过 → 安装依赖包（R \texttt{install.packages} / Python \texttt{pip install} / Stata \texttt{ssc install}）
  \item Step 3不过 → 重新拉取qoder-skills-library（git pull）或安装缺失Skill
\end{itemize}
\end{warningbox}

% ===========================================================
\section{20分类Skill全景安装与环境检查总表}

% ===========================================================

\textbf{【本章特别要求】} 本节是学生查阅频率最高的"环境准备速查表"，覆盖全部20个分类中的\textbf{11类高关联度Skill}（其余9类列出位置但默认低频）。每个Skill都包含"安装指令 + 环境检查命令"两部分。

\subsection{01-ideation 选题类（6个Skill）}
\textbf{idea-generation} / \texttt{research-ideation} / \texttt{ideation-spark}等：低依赖，只需Cursor/Claude Code调用。
\begin{lstlisting}[style=shell]
ls qoder-skills-library/01-ideation/   # 验证存在性
\end{lstlisting}
\textbf{brainstorming Skill}（全局可用）：无需额外环境，调用即可。

\subsection{02-literature 文献综述类（10个Skill）}
\textbf{find-literature} / \texttt{lit-review-assistant} / \texttt{academic-researcher} / \texttt{systematic-literature-review}：需要联网访问Semantic Scholar / arXiv / OpenAlex。
\begin{lstlisting}[style=shell]
curl -I https://api.semanticscholar.org/ | head -1   # 应返回 HTTP 200
curl -I https://export.arxiv.org/ | head -1          # arXiv API
\end{lstlisting}
\textbf{pdf-extraction Skill}：依赖Python pdfplumber。
\begin{lstlisting}[style=shell]
pip install pdfplumber
python3 -c "import pdfplumber; print(pdfplumber.__version__)"
\end{lstlisting}
\textbf{scansci-pdf Skill}：需要校园VPN配置（详见SKILL.md）。
\begin{lstlisting}[style=shell]
cat mcp.json | jq '.mcpServers."scansci-pdf"'
\end{lstlisting}

\subsection{04-data 数据获取与清洗类（3个Skill）—— 重点}
\textbf{tushare-finance Skill}：依赖Tushare Token。
\begin{lstlisting}[style=shell]
# 申请 Token：https://tushare.pro/register
export TUSHARE_TOKEN='your_token_here'
python3 -c "import tushare; print(tushare.__version__)"
\end{lstlisting}
\textbf{api-data-fetcher Skill}：依赖requests库 + API Key。\texttt{pip install requests pandas}

\textbf{worldbank-data360 / oecd-data}：依赖联网。
\begin{lstlisting}[style=shell]
curl -I https://api.worldbank.org/ | head -1
curl -I https://stats.oecd.org/ | head -1
\end{lstlisting}

\subsection{05-analysis 统计分析与仿真类（11个Skill）—— 最重点}
\begin{itemize}[leftmargin=1em]
  \item \textbf{stata Skill}：见§R/Stata环境（已详述）
  \item \textbf{stata-regression Skill}：stata依赖 + reghdfe/estout包
  \item \textbf{r-econometrics Skill}：见§R环境（已详述）
  \item \textbf{python-panel-data Skill}：见§Python环境 + linearmodels/PyFixest
  \item \textbf{data-analysis Skill}：见§Python环境
  \item \textbf{abm-analyzer / agent-based-modeling / agricultural-insurance-abm Skill}：依赖Mesa框架
\end{itemize}
\begin{lstlisting}[style=shell]
pip install mesa numpy pandas matplotlib
python3 -c "import mesa; print(mesa.__version__)"
\end{lstlisting}
\textbf{latex-tables Skill}：依赖esttab（R/Stata）或modelsummary（R）。
\begin{lstlisting}[style=shell]
Rscript -e 'library(modelsummary); packageVersion("modelsummary")'
stata-mp -e "ssc install estout, replace"
\end{lstlisting}

\subsection{06-experiment 实验流程类（5个Skill）}
\textbf{experiment-plan} / \texttt{experiment-tracker} / \texttt{ml-experiment-tracking}：低依赖，仅需Python。

\textbf{ab-test-analysis Skill}：依赖scipy/statsmodels。

\subsection{07-writing 论文写作类（7个Skill）—— 重点}
\textbf{paper-write} / \texttt{academic-paper-writer Skill}：依赖LaTeX。
\begin{lstlisting}[style=shell]
which xelatex      # macOS: /Library/TeX/texbin/xelatex 或 /usr/local/bin/xelatex
xelatex --version  # 应输出 XeTeX 3.141592653
\end{lstlisting}
\textbf{latex-tables Skill}：见05-analysis。

\textbf{humanizer-zh Skill}：依赖联网调用AI检测API。

\textbf{paper-compile Skill}：依赖完整LaTeX环境 + latexmk。
\begin{lstlisting}[style=shell]
which latexmk && latexmk --version
\end{lstlisting}

\subsection{09-review 审稿类（11个Skill）}
\textbf{paper-review} / \texttt{peer-review} / \texttt{research-review} / \texttt{research-critique} / \texttt{scientific-critical-thinking} / \texttt{scholar-evaluation}：低依赖（多为LLM推理）。

\textbf{auto-review-loop} / \texttt{auto-review-loop-llm} / \texttt{auto-review-loop-minimax Skill}：依赖OpenAI/Codex/MiniMax API Key。
\begin{lstlisting}[style=shell]
export OPENAI_API_KEY='sk-...'   # auto-review-loop / auto-review-loop-llm
export MINIMAX_API_KEY='...'     # auto-review-loop-minimax
\end{lstlisting}
\textbf{auto-paper-improvement-loop Skill}：依赖GPT-5.4访问。
\begin{lstlisting}[style=shell]
export GPT54_API_KEY='sk-...'
\end{lstlisting}

\subsection{10-communication 演示展示类（7个Skill）—— 重点}
\begin{itemize}[leftmargin=1em]
  \item \textbf{beamer-presentation / paper-slides / paper-poster Skill}：见§运行环境准备-LaTeX Beamer
  \item \textbf{pptx Skill}：见§Node.js + LibreOffice环境
  \item \textbf{qoderwork-ppt Skill}：见§Puppeteer + LibreOffice环境
  \item \textbf{ppt-master Skill}：AI驱动多格式SVG生成，需联网 + Gemini/Claude API
\end{itemize}
\begin{lstlisting}[style=shell]
export GEMINI_API_KEY='your_key'    # https://aistudio.google.com/app/apikey
export ANTHROPIC_API_KEY='sk-ant-...'
node --version && which libreoffice && which pdftoppm
\end{lstlisting}
\textbf{paper-figure / econ-visualization / paper-illustration Skill}：见05-analysis与§R环境。

\textbf{mermaid-diagram Skill}：依赖mermaid CLI。
\begin{lstlisting}[style=shell]
npm install -g @mermaid-js/mermaid-cli
mmdc --version
\end{lstlisting}

\subsection{11-perspective 思维顾问类（2个Skill）}
\textbf{china-finance-knowledge-system-v6 Skill}：低依赖（本地知识库）。

\textbf{munger-thinking / taleb-thinking}（如存在）：低依赖。

\subsection{12-finance 金融与商业类（11个Skill）—— 最重点}
\textbf{finance-expert} / \texttt{financial-modeling} / \texttt{backtesting-trading-strategies Skill}：低依赖（多为LLM推理）。

\textbf{stock-analysis / china-stock-analysis Skill}：见§Tushare Token + china-stock-mcp MCP。

\textbf{yahoo-finance Skill}：依赖yfinance库。
\begin{lstlisting}[style=shell]
pip install yfinance
python3 -c "import yfinance; print(yfinance.__version__)"
\end{lstlisting}
\textbf{openbb-finance Skill}：依赖OpenBB Terminal。
\begin{lstlisting}[style=shell]
pip install openbb
openbb-build
\end{lstlisting}
\textbf{financial-unit-economics / saas-economics-efficiency-metrics / startup-financial-modeling / questionnaire-design Skill}：低依赖。

\subsection{16-office 文档办公类（8个Skill）—— 重点}
\textbf{docx Skill}：依赖python-docx。
\begin{lstlisting}[style=shell]
pip install python-docx
python3 -c "import docx; print(docx.__version__)"
\end{lstlisting}
\textbf{xlsx Skill}：依赖openpyxl/pandas。\texttt{pip install openpyxl pandas}

\textbf{pdf Skill}：依赖pypdf/pdfplumber。\texttt{pip install pypdf pdfplumber}

\textbf{cli-anything / pdf-office-suite / doc-suite / spreadsheet-suite Skill}：依项检查。

\textbf{pptx Skill}：见10-communication。

\textbf{office-powerpoint / office-word MCP}：需mcp.json中启用。

\subsection{19-devtools 元工具类（20个Skill）—— 本章重点}
详见§元Skill体系章节（\texttt{create-skill} / \texttt{find-skills} / \texttt{skill-auditor} / \texttt{write-skill-readme} / \texttt{install-skill-dependency} / \texttt{luban}）。

其他工具：\texttt{github-cli} / \texttt{claude-code} / \texttt{cursor-cli}等。
\begin{lstlisting}[style=shell]
gh --version       # GitHub CLI
claude --version   # Claude Code（可选）
cursor --version   # Cursor IDE（可选）
\end{lstlisting}

\subsection{环境速查命令（一键检查所有11类高关联度Skill）}

复制以下脚本到本地运行，5秒内输出所有Skill的可用性报告。

\begin{lstlisting}[style=shell]
cat > /tmp/check_all_skills.sh << 'EOF'
#!/bin/bash
echo "=== 环境检查报告 ==="
echo "[R]    $(R --version 2>/dev/null | head -1 || echo MISSING)"
echo "[Stata] $(stata-mp --version 2>/dev/null | head -1 || echo MISSING)"
echo "[Py]    $(python3 --version 2>/dev/null || echo MISSING)"
echo "[Node]  $(node --version 2>/dev/null || echo MISSING)"
echo "[LaTeX] $(xelatex --version 2>/dev/null | head -1 || echo MISSING)"
echo "[LibreOffice] $(libreoffice --version 2>/dev/null | head -1 || echo MISSING)"
echo "[poppler] $(pdftoppm -v 2>&1 | head -1 || echo MISSING)"
echo "[mermaid] $(mmdc --version 2>/dev/null || echo MISSING)"
echo "[GitHub] $(curl -s -o /dev/null -w '%{http_code}' https://github.com)"
echo "[Tushare] $([ -n "$TUSHARE_TOKEN" ] && echo SET || echo MISSING)"
echo "[Gemini] $([ -n "$GEMINI_API_KEY" ] && echo SET || echo MISSING)"
echo "=== Skill库存盘点 ==="
ls qoder-skills-library/ | head -30
EOF
chmod +x /tmp/check_all_skills.sh && /tmp/check_all_skills.sh
\end{lstlisting}

% ===========================================================
\section{金融Skill生态全景}
% ===========================================================

\subsection{金融Skill分类概览}

以下是本课程技能库中与银行金融直接相关的Skill（原表13行，现扩充为\textbf{5大类共30+个Skill}）：

\begin{table}[H]
\centering
\begin{tabular}{lp{2.8cm}p{8cm}}
\toprule
\textbf{大类} & \textbf{Skill名称} & \textbf{核心能力} \\
\midrule
\multirow{6}{*}{\textbf{数据获取类}} & \texttt{tushare-finance} & tushare Pro API获取中国金融市场数据（股票/基金/期货） \\
& \texttt{china-stock-analysis} & A股数据获取 + 财务指标计算 \\
& \texttt{yahoo-finance} & 美股/港股全球股票数据（yfinance库） \\
& \texttt{openbb-finance} & OpenBB Terminal多资产多市场金融数据 \\
& \texttt{api-data-fetcher} & 通用API数据拉取框架（多数据源适配） \\
& \texttt{finance-news} & 实时财经新闻抓取 + 情绪分析 \\
\midrule
\multirow{8}{*}{\textbf{专业分析类}} & \texttt{finance-expert} & DCF估值、信用风险评估、DuPont分析、五C模型 \\
& \texttt{financial-modeling} & 构建DCF/LBO/M\&A模型 + 敏感性分析 \\
& \texttt{backtesting-trading-strategies} & 量化策略回测 + Sharpe/回撤计算 \\
& \texttt{financial-unit-economics} & 单位经济学分析（SaaS/金融科技） \\
& \texttt{stock-analysis} & 全球股票分析（含Yahoo Finance数据） \\
& \texttt{saas-economics-efficiency-metrics} & SaaS企业效率指标分析 \\
& \texttt{startup-financial-modeling} & 初创企业财务建模 \\
& \texttt{questionnaire-design} & 学术调查问卷设计（金融普惠等方向） \\
\midrule
\multirow{8}{*}{\textbf{计量与数据科学类}} & \texttt{stata} & Stata全面参考（.do文件/计量/图表/包） \\
& \texttt{stata-regression} & reghdfe/estout一键输出三线回归表 \\
& \texttt{r-econometrics} & R fixest包（高维FE/IV/DiD）+ 计量生态 \\
& \texttt{python-panel-data} & Python面板数据回归（linearmodels/PyFixest） \\
& \texttt{data-analysis} & 通用数据分析与预处理 \\
& \texttt{abm-analyzer} & Agent-Based Modeling仿真分析 \\
& \texttt{agent-based-modeling} & ABM建模与仿真 \\
& \texttt{agricultural-insurance-abm} & 农业保险ABM专题仿真 \\
\midrule
\multirow{8}{*}{\textbf{文档与可视化类}} & \texttt{xlsx} & Excel读写 + 财务模型 + 公式 + 图表 \\
& \texttt{docx} & Word报告生成（封面/目录/页码） \\
& \texttt{pptx} & PPT演示文稿生成与编辑（python-pptx） \\
& \texttt{pdf} & PDF读写 + 表格提取 \\
& \texttt{qoderwork-ppt} & 14种HTML模板 + 一键生成PPTX \\
& \texttt{cli-anything} & 通用CLI命令适配框架 \\
& \texttt{econ-visualization} & 学术发表级图表（R ggplot2） \\
& \texttt{paper-figure} & 论文图表（Python matplotlib） \\
& \texttt{mermaid-diagrams} & Mermaid流程图/时序图/甘特图 \\
\midrule
\multirow{5}{*}{\textbf{思维与方法论类}} & \texttt{consulting-frameworks} & McKinsey 7S / BCG / 波特五力 / SWOT \\
& \texttt{humanizer-zh} & AI文本检测与中文人文改写 \\
& \texttt{china-finance-knowledge-system-v6} & 中国金融学知识体系 \\
& \texttt{systematic-literature-review} & 系统性文献综述生成 \\
& \texttt{lit-review-assistant} & 文献综述辅助助手 \\
\bottomrule
\end{tabular}
\end{table}

\begin{notebox}[Skill分类维度]
上述Skill可按两个维度理解：
\begin{itemize}[leftmargin=1em]
  \item \textbf{按功能定位}：数据获取类（获取原始信息）、分析类（产生专业判断）、文档类（输出标准成果物）、咨询类（提供框架思维）
  \item \textbf{按银行业务场景}：零售业务（客户分群/产品推荐）、公司业务（信用评估/贷款定价）、风险管理（舆情监控/风险预警）、运营管理（报告生成/数据报表）
\end{itemize}
\end{notebox}

\subsection{Skill市场与安装方式}

Skill的获取主要有三种途径：

\subsubsection{skills.sh市场}

skills.sh是当前最活跃的Skill分发平台（\url{https://skills.sh}），类似于npm之于Node.js生态。在该平台上可以浏览、搜索和安装各类Skill。

\begin{lstlisting}[style=shell]
# 搜索金融类 Skill
npx skills search finance

# 查看某个 Skill 的详细信息
npx skills info finance-expert
\end{lstlisting}

\subsubsection{npm安装方式}

每个Skill都通过npm进行分发，安装命令格式为\texttt{npx skills add <组织>/<仓库>@<技能名>}：

\begin{lstlisting}[style=shell]
# 标准安装格式
npx skills add <org>/<repo>@<skill-name> -y

# 示例
npx skills add anthropics/skills@docx -y
npx skills add personamanagmentlayer/pcl@finance-expert -y
\end{lstlisting}

\subsubsection{本地Skill管理}

对于自编写的Skill，只需将其放置在项目的\texttt{.agents/skills/}目录下即可被AI IDE自动发现：

\begin{lstlisting}[style=shell]
# 本地 Skill 目录结构
project/
├── .agents/
│   └── skills/
│       └── bank-risk-alert/
│           └── SKILL.md    # 自编写的Skill文件
├── reports/
└── ...
\end{lstlisting}

\subsection{环境准备}

在开始本章实验之前，请确认以下环境要求已满足：

\begin{enumerate}[leftmargin=2em]
  \item TRAE / Qoder 已更新到最新版本
  \item 实验二MCP已配通（特别是excel、word、ppt MCP）
  \item 安装本实验所需的全部Skill（在终端执行以下命令）：
\end{enumerate}

\begin{lstlisting}[style=shell]
# 安装 Skills CLI（如未安装）
npm install skills -g

# 一键安装本实验所需的 6 个 Skill
npx skills add anthropics/skills@docx -y
npx skills add personamanagmentlayer/pcl@finance-expert -y
npx skills add anthropics/skills@xlsx -y
npx skills add softaworks/agent-toolkit@mermaid-diagrams -y
npx skills add sundial-org/awesome-openclaw-skills@finance-news -y
npx skills add aznatkoiny/zai-skills@consulting-frameworks -y

# 验证安装结果
npx skills ls
\end{lstlisting}

\begin{notebox}[安装说明]
\begin{itemize}[leftmargin=1em]
  \item \texttt{npx skills add} 会自动从Skills市场（\url{https://skills.sh}）下载并安装到当前项目
  \item 安装完成后可用 \texttt{npx skills ls} 查看已安装的技能列表
  \item 如安装失败，检查网络连接或使用 \texttt{npm config set registry https://registry.npmmirror.com} 切换镜像
  \item TRAE用户：也可在Builder模式中直接告诉AI「请帮我安装docx / xlsx / finance-expert技能」
\end{itemize}
\end{notebox}

\textbf{验证Skill系统可用}：在对话框输入以下提示词

\begin{lstlisting}
请用 mermaid-diagram skill 画一个简单流程图：开始→处理→结束
\end{lstlisting}

如果AI能生成.mmd文件 = Skill系统正常。

Skill存放位置（了解即可）：Qoder = \texttt{\%APPDATA\%\textbackslash Qoder\textbackslash skills\textbackslash}；TRAE内置无需手动管理。

% ===========================================================
\section{元Skill体系：管理Skill的Skill}

% ===========================================================

元Skill（Meta-Skill）是\textbf{管理Skill的Skill}，是本章的亮点。本节覆盖6个核心元Skill，按"创建→发现→审计→打磨→发布"的生命周期组织。

\subsection{create-skill（创建元Skill）}
\textbf{触发词}："创建一个新Skill"、"为XX场景生成Skill"、"按SKILL.md模板生成"。

\textbf{能力清单}：跟随交互向导生成完整SKILL.md（含frontmatter、argument-hint、allowed-tools、workflow\_stage、tags）。

\textbf{环境检查}：依赖Node.js + Cursor/Claude Code（其他IDE可能不支持完整向导）。
\begin{lstlisting}[style=shell]
node --version   # >= 18.x
npm --version
ls qoder-skills-library/19-devtools/create-skill/SKILL.md
\end{lstlisting}

\subsection{find-skills（发现元Skill）}
\textbf{触发词}："找Skill"、"有什么Skill能做XX"、"skills.sh市场"、"Skill发现"。

\textbf{能力清单}：从skills.sh生态、GitHub anthropics/skills仓库、本地qoder-skills-library、npm registry发现Skill。

\textbf{环境检查}：依赖网络访问https://skills.sh与https://github.com。
\begin{lstlisting}[style=shell]
curl -I https://skills.sh/  | head -1   # 应返回 HTTP 200
ping -c 1 github.com                    # 验证 GitHub 可达
\end{lstlisting}

\subsection{skill-auditor（审计元Skill）—— 五维审计自动化}
\textbf{触发词}："审计这个Skill"、"检查Skill质量"、"SKILL.md五维评分"。

\textbf{能力清单}：对任意SKILL.md进行五维自动评分（触发词准确性、输出完整性、数据依赖、合规性、改进建议），输出结构化审计报告。

\textbf{环境检查}：独立运行，无需额外依赖。
\begin{lstlisting}[style=shell]
ls qoder-skills-library/19-devtools/skill-auditor/SKILL.md
qoder skill run skill-auditor --target=.agents/skills/bank-risk-alert/SKILL.md
\end{lstlisting}

\subsection{write-skill-readme（文档生成元Skill）}
\textbf{触发词}："生成Skill README"、"补README"、"Skill文档"。

\textbf{能力清单}：根据SKILL.md + 示例自动生成README.md（含徽章、使用场景、示例代码）。

\subsection{install-skill-dependency（环境诊断元Skill）}
\textbf{触发词}："Skill环境诊断"、"检查依赖"、"Skill装不起来"。

\textbf{能力清单}：诊断Skill调用前的所有环境依赖（R包/Stata外部命令/Python包/Node模块/API Key），输出缺失清单与一键安装脚本。

\textbf{环境检查}：本元Skill本身依赖Python 3.8+。
\begin{lstlisting}[style=shell]
python3 --version   # >= 3.8
pip show skill-diagnostic 2>/dev/null || pip install skill-diagnostic
\end{lstlisting}

\subsection{luban（打磨元Skill）}
\textbf{触发词}："打磨Skill"、"Skill升级"、"鲁班打磨"、"Skill产品化"。

\textbf{能力清单}：工匠式五动作（验料/访行/过尺/慢刨/回炉），产出《Skill打磨报告》 + 可替换改写片段 + 出师证书。

\subsection{元Skill协同工作流（实战提示词）}

\begin{lstlisting}
我要为本行创建一个"银行反洗钱可疑交易报告生成Skill"（名为 aml-suspicious-report），请按以下流程协助：
1. 调用 create-skill 生成 SKILL.md 初始草稿
2. 调用 find-skills 查找生态中同类Skill（如 finance-expert、bank-risk-alert）作为参考
3. 调用 skill-auditor 对生成的草稿进行五维自动审计，输出改进点
4. 按 skill-auditor 的反馈迭代修改（最多 2 轮）
5. 调用 write-skill-readme 生成 README.md
6. 调用 install-skill-dependency 检查Skill是否能跑起来
7. （可选）调用 luban 进一步打磨该Skill，准备发布到 qoder-skills-library
\end{lstlisting}

\textbf{预期产出}：
\begin{itemize}[leftmargin=1em]
  \item \texttt{aml-suspicious-report/SKILL.md}（草稿定稿）
  \item \texttt{aml-suspicious-report/README.md}（自动生成）
  \item \texttt{aml-suspicious-report/audit\_report.md}（五维审计报告）
  \item \texttt{aml-suspicious-report/env\_check.txt}（环境检查日志）
\end{itemize}

% ===========================================================
\section{Skill调用实战}
% ===========================================================

本节通过三个由浅入深的实战任务，带你掌握金融Skill的调用方法。每个任务对应一个真实的银行业务场景。

\begin{table}[H]
\centering
\begin{tabular}{lp{5.5cm}l}
\toprule
\textbf{时段} & \textbf{内容} & \textbf{交付物} \\
\midrule
0--20 min & Skill概念 + 金融Skill全景介绍 & --- \\
20--70 min & \textbf{任务1}：docx Skill客户回访报告 & .docx文件 \\
70--120 min & \textbf{任务2}：finance-expert信用风险分析 & 风险分析报告 \\
120--160 min & \textbf{任务3}：xlsx Skill贷款定价模型 & .xlsx文件 \\
160--195 min & \textbf{任务4}：自写bank-risk-alert Skill & SKILL.md \\
195--225 min & \textbf{任务5}：Skill安全审计 & 审计报告 \\
225--240 min & 总结 + 学生demo互看 & --- \\
\bottomrule
\end{tabular}
\end{table}

% -----------------------------------------------------------
\subsection{docx Skill：银行客户回访报告（\ding{72}）}
% -----------------------------------------------------------

\textbf{场景}：你是某城商行客户经理，需要为VIP客户C00001（张**）生成一份规范的电话回访报告。

\textbf{技能安装}（如环境准备阶段未安装）：

\begin{lstlisting}[style=shell]
npx skills add anthropics/skills@docx -y
\end{lstlisting}

\textbf{操作提示词}：

\begin{lstlisting}
请执行以下操作以完成客户电话回访报告的生成：

1. 安装 docx 技能环境：
   - 前往热门技能市场搜索并安装"docx"技能包
   - 通过 npm 全局安装 docx-js 工具：执行命令 "npm install -g docx"
   - 验证 pandoc 与 LibreOffice 软件是否已正确安装（用于后续 PDF 格式校验）
   - 完成安装后回复确认信息："docx skill 已就绪"

2. 使用已安装的 docx 技能为客户 C00001（张**）生成《客户电话回访报告》：
   - 创建结构完整的 Word 文档，包含以下 6 个标准章节：
     * 客户基本信息（客户编号 C00001、姓名、资产等级、开户行、客户经理）
     * 回访背景（本次电话回访的目标：定期存款到期提醒 + 理财产品推荐）
     * 沟通要点（详细记录电话沟通中的关键内容与讨论事项）
     * 客户反馈（客户对现有服务的满意度评分 1-5 + 具体意见）
     * 产品意向（客户对大额存单/结构性存款/基金定投的兴趣程度）
     * 跟进计划（含时间节点、责任人、后续行动项）
   - 文档格式要求：
     * 添加正式封面，包含银行 Logo 位置、报告标题、日期
     * 生成自动更新的目录
     * 添加规范的页码，页眉含「内部文件·密级：内部」
   - 保存路径："reports/客户回访报告_C00001.docx"

3. 完成后请确认文档已成功生成并符合上述所有要求。
\end{lstlisting}

\textbf{LibreOffice安装提示词}（如需要使用）：

\begin{lstlisting}
请帮我检查本机是否安装了 LibreOffice。当前系统为 Windows，shell 为 PowerShell，请按以下流程执行：

【第一步：探测安装状态】
依次执行下列检测，任一命中即视为"已安装"：
1. 检测默认安装目录是否存在 soffice.exe
2. 检测 PATH 中是否有 soffice
3. 检测注册表卸载项

【第二步：未安装则自动下载并静默安装】
方案 A（首选，winget）：
winget install --id TheDocumentFoundation.LibreOffice -e --silent --accept-source-agreements --accept-package-agreements

方案 B（备选，Chocolatey）：
choco install libreoffice-fresh -y

方案 C（兜底，官方 MSI 直链下载并静默安装）

【第三步：安装后验证】
1. 重新检测 soffice 是否存在
2. 执行 soffice --version 打印版本
3. 提示用户：如需在 PowerShell 中直接使用 soffice 命令，可将安装目录加入用户 PATH
\end{lstlisting}

\begin{theorybox}[银行文档规范要求]
在银行业务中，正式文档必须遵循严格的格式规范。使用docx Skill生成文档时，应特别注意以下要素：

\begin{enumerate}[leftmargin=1em]
  \item \textbf{密级标记}：根据《商业银行数据安全管理办法》，银行文档应在页眉或封面明确标注密级——"公开""内部""机密""绝密"。客户回访报告属于"内部"级别。
  \item \textbf{文号规范}：正式行文应包含发文字号，格式为"〔年份〕序号"，如"XX银行客户部〔2026〕015号"。
  \item \textbf{签章位置}：重要文档需预留签章区域，包括"编制人""复核人""审批人"的签名栏和日期栏。
  \item \textbf{版本控制}：文档封面应标注版本号和修订日期，如"V1.0 / 2026-06-08"。
  \item \textbf{页眉页脚}：页眉通常包含银行名称和密级标记，页脚包含页码和文档编号。
\end{enumerate}

在提示词中明确这些规范要求，可以有效提升AI生成文档的合规性。
\end{theorybox}

% -----------------------------------------------------------
\subsection{finance-expert Skill：信用风险评估（\ding{72}\ding{72}）}
% -----------------------------------------------------------

\textbf{场景}：银行风控部门需要对一家中小企业客户进行信用风险评估，辅助审批500万元流动资金贷款。

\textbf{技能安装}（如环境准备阶段未安装）：

\begin{lstlisting}[style=shell]
npx skills add personamanagmentlayer/pcl@finance-expert -y
\end{lstlisting}

\begin{theorybox}[信用风险评估的五C模型]
五C模型（Five C's of Credit）是商业银行信用风险评估的经典框架，由五个以C开头的维度构成：

\begin{enumerate}[leftmargin=1em]
  \item \textbf{Character}（品格）：借款人的还款意愿与信用历史。考察其过往贷款记录、违约历史、诉讼情况。在中小企业评估中，实际控制人的个人信用同样关键。
  \item \textbf{Capacity}（能力）：借款人的还款能力。核心指标包括资产负债率、流动比率、利息保障倍数、经营性现金流与贷款本息的比值等。
  \item \textbf{Capital}（资本）：借款人的自有资金实力。资本充足的企业抗风险能力更强，银行也倾向于看到借款人投入了足够的自有资金。
  \item \textbf{Collateral}（抵押）：用于担保贷款的资产质量。抵押物的评估价值、变现能力、抵押率（LTV）都是重要考量。
  \item \textbf{Conditions}（条件）：宏观经济与行业环境。即使个体财务健康，行业下行周期也可能导致系统性风险。
\end{enumerate}

在AI辅助信用评估中，finance-expert Skill主要聚焦于\textbf{Capacity}和\textbf{Capital}两个维度的量化分析，而\textbf{Character}、\textbf{Collateral}和\textbf{Conditions}则需要结合外部数据与人工判断。
\end{theorybox}

\textbf{操作提示词}：

\begin{lstlisting}
请使用 finance-expert 技能完成以下银行信用风险分析任务：

【背景】
某城商行收到客户"恒达机械制造有限公司"的贷款申请：
- 申请额度：500 万元流动资金贷款
- 期限：12 个月
- 用途：原材料采购

【第 1 步】构建企业信用评估框架
请基于以下模拟财务数据建立评估模型：
- 近三年营收：2023年 3200万 / 2024年 3800万 / 2025年 4100万
- 近三年净利润：180万 / 220万 / 250万
- 资产总额：2800万，负债总额：1600万
- 应收账款周转天数：85天
- 存货周转天数：60天
- 经营性现金流：近12个月净流入 320万

【第 2 步】计算关键风控指标
1. 资产负债率、流动比率、速动比率
2. Altman Z-Score（适配中国制造业版本）
3. 利息保障倍数（假设年利率 5.5%）
4. 贷款偿还能力系数 = 经营性现金流 / 贷款本息

【第 3 步】生成风险评估报告
输出结构化报告（Markdown 格式），包含：
1. 企业画像概要
2. 财务健康度评分（百分制）
3. 风险等级判定（低/中/高/极高）
4. 审批建议（批准/有条件批准/拒绝）
5. 风险缓释措施建议（如需要抵押/担保/降额等）

保存到 reports/credit_risk_hengda.md
\end{lstlisting}

\textbf{学生练习}：修改企业财务数据（如资产负债率提高到75\%），观察风险评级如何变化。

\textbf{拓展A：银行同业对标分析}

\begin{lstlisting}
请使用 finance-expert 技能完成银行同业对标分析：

对比分析以下 4 家银行的经营能力（使用模拟数据或公开年报数据）：
1. 工商银行（大型国有银行代表）
2. 招商银行（股份制银行代表）
3. 北京银行（城商行代表）
4. 常熟银行（农商行代表）

对比维度：
- 盈利能力：ROE / ROA / 净息差（NIM）
- 资产质量：不良贷款率 / 拨备覆盖率 / 关注类占比
- 资本充足：核心一级资本充足率 / 资本充足率
- 成长性：营收增速 / 净利润增速 / 贷款增速

输出：对比表格 + 雷达图描述 + 200 字总结
保存到 reports/bank_peer_comparison.md
\end{lstlisting}

% -----------------------------------------------------------
\subsection{xlsx Skill：贷款FTP定价模型（\ding{72}\ding{72}\ding{72}）}
% -----------------------------------------------------------

\textbf{场景}：银行资产负债部需要构建一个贷款FTP定价模型，用于指导客户经理合理报价。

\textbf{技能安装}（如环境准备阶段未安装）：

\begin{lstlisting}[style=shell]
npx skills add anthropics/skills@xlsx -y
\end{lstlisting}

\begin{theorybox}[FTP转移定价理论基础]
\textbf{内部资金转移定价}（Funds Transfer Pricing，FTP）是商业银行内部经营管理的关键工具。其核心思想是：银行将资金集中到"资金中心"（类似内部银行），各业务条线从资金中心"买入"资金用于放贷，或"卖出"资金至资金中心。

FTP定价的基本公式为：
\[
\text{贷款利率} = \text{FTP资金成本} + \text{风险溢价} + \text{运营成本} + \text{资本占用成本} + \text{目标利润}
\]

各组成部分的含义：
\begin{itemize}[leftmargin=1em]
  \item \textbf{FTP资金成本}：资金中心向业务条线收取的资金价格，通常参考同期限Shibor或内部收益率曲线。一年期FTP约2.8\%，三年期约3.2\%。
  \item \textbf{风险溢价}：根据客户信用等级差异化定价，AAA级仅0.3\%，BB级则高达3.0\%，反映预期损失（EL = PD $\times$ LGD $\times$ EAD）。
  \item \textbf{运营成本}：人工、系统、网点等分摊成本，通常为0.3\%--0.5\%。
  \item \textbf{资本占用成本}：监管要求银行持有一定比例的资本金以覆盖非预期损失，公式为：贷款金额 $\times$ 风险权重 $\times$ 资本充足率 $\times$ 资本回报率。
  \item \textbf{目标利润}：银行股东要求的最低回报，通常为0.1\%--0.3\%。
\end{itemize}

FTP定价的意义在于：将利率风险从业务条线转移至资金中心集中管理，使客户经理的定价决策基于完全成本而非边际成本，避免"赔本赚吆喝"。
\end{theorybox}

\textbf{操作提示词}：

\begin{lstlisting}
请使用 xlsx 技能创建银行贷款定价模型 Excel 文件：

【文件名】reports/loan_pricing_model.xlsx

【Sheet 1：定价参数】
构建贷款定价公式：贷款利率 = 资金成本 + 风险溢价 + 运营成本 + 资本占用成本 + 目标利润

参数设置（使用 Excel 公式，不硬编码计算结果）：
- 资金成本（FTP）：一年期 2.8%，三年期 3.2%
- 风险溢价：按信用等级差异化
  * AAA级：0.3%  AA级：0.6%  A级：1.0%  BBB级：1.8%  BB级：3.0%
- 运营成本：0.4%（含人工、系统、网点分摊）
- 资本占用成本：贷款金额 × 风险权重(100%) × 资本充足率(12.5%) × 资本回报率(12%)
- 目标利润：0.2%

【Sheet 2：客户定价测算】
为 5 位不同客户自动计算贷款报价：
| 客户名 | 金额(万) | 期限 | 信用等级 | 抵押方式 | 最终利率 |
| 恒达机械 | 500 | 1年 | A | 厂房抵押 | =公式计算 |
| 绿源农业 | 200 | 3年 | BBB | 信用 | =公式计算 |
| 华新科技 | 1000 | 1年 | AA | 应收质押 | =公式计算 |
| 鑫达贸易 | 300 | 1年 | BB | 保证担保 | =公式计算 |
| 瑞丰地产 | 2000 | 3年 | A | 土地抵押 | =公式计算 |

【Sheet 3：敏感性分析】
- 当 FTP 上浮/下浮 20bp 时，各客户利率如何变化
- 当信用等级上调/下调一级时，利率如何变化
- 用条件格式标注：利率 < 4% 绿色，4-6% 黄色，> 6% 红色

【格式要求】
- 所有计算必须使用 Excel 公式（=SUM, =VLOOKUP, =IF 等）
- 表头加粗 + 蓝色背景
- 百分比统一保留 2 位小数
- 金额使用千位分隔符
\end{lstlisting}

\textbf{学生练习}：尝试添加「抵押折扣」字段——有足值抵押的客户可享受风险溢价8折优惠。

\textbf{拓展B：银行客户分群RFM模型}

\begin{lstlisting}
请使用 xlsx 技能构建银行客户 RFM 分群模型：

1. 在 Sheet1 创建 30 条模拟客户交易数据：
   列：客户ID / 姓名 / 最近交易日期 / 近12月交易频次 / 近12月交易总额

2. 在 Sheet2 计算 RFM 分值：
   - R（Recency）：最近交易距今天数，越少越好，1-5分
   - F（Frequency）：交易频次，越多越好，1-5分
   - M（Monetary）：交易金额，越大越好，1-5分
   - RFM总分 = R + F + M

3. 在 Sheet3 客户分群：
   - 重要价值客户（RFM >= 12）
   - 重要发展客户（R高 F低 M高）
   - 重要保持客户（R低 F高 M高）
   - 一般客户（RFM < 8）

4. 输出各分群客户数、平均资产、营销建议
\end{lstlisting}

% -----------------------------------------------------------
\subsection{tushare-finance + china-stock-analysis Skill：银行股估值对比（\ding{72}\ding{72}\ding{72}）}
% -----------------------------------------------------------

\textbf{场景}：银行业研究员需要对比 A 股银行股与 H 股银行股的估值水平，为本行投资决策提供数据支持。

\textbf{技能安装}：

\begin{lstlisting}[style=shell]
# Tushare Skill 前置
pip install tushare pandas
export TUSHARE_TOKEN='your_token_here'    # 申请：https://tushare.pro/register
# 验证
python3 -c "import tushare; pro=tushare.pro_api(); print(pro.stock_basic(list_status='L', limit=1))"
\end{lstlisting}

\textbf{3 步环境检查}：
\begin{lstlisting}[style=shell]
# Step 1：路径检查
which python3
# Step 2：依赖检查
python3 -c "import tushare; print(tushare.__version__)"
[ -n "$TUSHARE_TOKEN" ] && echo "TUSHARE_TOKEN SET" || echo "MISSING"
# Step 3：Skill 可访问性
ls qoder-skills-library/04-data/tushare-finance/SKILL.md
\end{lstlisting}

\textbf{操作提示词（Skill 化包装）}：

\begin{lstlisting}
请调用 tushare-finance 技能完成 A 股银行股估值对比分析：

【调用 Skill】tushare-finance、china-stock-analysis
【背景】为本行投资决策提供 A/H 股银行股估值对比
【输入数据】
  - A 股银行：工商银行(601398)、招商银行(600036)、北京银行(601169)、常熟银行(601128)
  - H 股银行：工商银行(1398.HK)、招商银行(3968.HK)
  - 时间范围：2010-01-01 至 2025-12-31
【输出要求】
  - reports/bank_valuation_compare.xlsx（A/H 股 PE/PB 对比表）
  - figures/bank_pe_compare.pdf（PE 对比折线图）
  - reports/bank_valuation_analysis.md（分析报告 ≥800 字）
【环境检查】运行 /tmp/check_all_skills.sh 确认 Tushare Token 已配置
【原代码补充】
  import tushare as ts
  pro = ts.pro_api()
  df_a = pro.daily_basic(ts_code='601398.SH', start='20240101')
  df_h = pro.hk_daily(ts_code='01398.HK', start='20240101')
【失败回退】如 china-stock-mcp 不可用，直接使用 tushare 单独拉取
【预期结论】A 股银行 PE 较 H 股溢价率、H 股银行 PB 折价幅度、估值修复空间
\end{lstlisting}

\textbf{学生练习}：将样本银行扩展至 16 家全国性商业银行，绘制 A/H 股溢价率分布图。

% -----------------------------------------------------------
\subsection{backtesting-trading-strategies Skill：银行股量化回测（\ding{72}\ding{72}\ding{72}\ding{72}）}
% -----------------------------------------------------------

\textbf{场景}：银行金融市场部需要回测"双均线 + RSI + 海龟"三类经典策略在银行股上的表现，为自营交易提供量化依据。

\textbf{技能安装}：

\begin{lstlisting}[style=shell]
pip install backtrader vectorbt pandas numpy matplotlib
ls qoder-skills-library/12-finance/backtesting-trading-strategies/SKILL.md
\end{lstlisting}

\textbf{3 步环境检查}：
\begin{lstlisting}[style=shell]
which python3
python3 -c "import backtrader; print('backtrader OK')"
python3 -c "import vectorbt; print('vectorbt OK')" || echo "vectorbt MISSING (optional)"
ls qoder-skills-library/12-finance/backtesting-trading-strategies/SKILL.md
\end{lstlisting}

\textbf{操作提示词（Skill 化包装）}：

\begin{lstlisting}
请调用 backtesting-trading-strategies 技能完成银行股量化策略回测：

【调用 Skill】backtesting-trading-strategies、tushare-finance、econ-visualization
【背景】测试三种经典技术策略在 5 家银行股上的表现
【输入数据】
  - 标的：工商银行、招商银行、兴业银行、平安银行、宁波银行（2015-2025）
  - 策略：
    * 策略 1：双均线（MA5/MA20 金叉死叉）
    * 策略 2：RSI（超买超卖阈值 30/70）
    * 策略 3：海龟交易法（20 日突破）
  - 初始资金：100 万元 / 手续费 0.1% / 滑点 0.05%
【输出要求】
  - reports/strategy_backtest_results.xlsx（绩效对比表）
  - figures/backtest_curves.pdf（三策略净值曲线对比）
  - reports/backtest_report.md（分析报告，含夏普/最大回撤/年化收益）
【环境检查】运行 /tmp/check_all_skills.sh
【原代码补充】backtrader Strategy 类示例代码（见 SKILL.md）
【失败回退】如 backtesting-trading-strategies Skill 不可用，手写 backtrader 策略代码
【预期结论】哪种策略在银行股上最稳健？年化夏普比率对比
\end{lstlisting}

\textbf{学生练习}：加入止损规则（亏损 5\% 自动平仓），观察回测结果变化。

% -----------------------------------------------------------
\subsection{econ-visualization Skill：银行 ROE 雷达图（\ding{72}\ding{72}\ding{72}）}
% -----------------------------------------------------------

\textbf{场景}：银行年报披露期，需要为投资者关系部门快速生成各家银行的多维财务对比图。

\textbf{技能安装}（前置 R 环境准备见 §R 语言环境）：

\begin{lstlisting}[style=shell]
Rscript -e 'install.packages(c("fmsb", "ggplot2", "tidyr", "dplyr", "scales"))'
ls qoder-skills-library/05-analysis/econ-visualization/SKILL.md
\end{lstlisting}

\textbf{3 步环境检查}：
\begin{lstlisting}[style=shell]
which R
Rscript -e 'library(fmsb); library(ggplot2); cat("OK\n")'
ls qoder-skills-library/05-analysis/econ-visualization/SKILL.md
\end{lstlisting}

\textbf{操作提示词（Skill 化包装）}：

\begin{lstlisting}
请调用 econ-visualization 技能生成 5 家银行 ROE 多维雷达图：

【调用 Skill】econ-visualization、tushare-finance
【背景】投资者关系部门年度披露材料
【输入数据】tushare 拉取 5 家银行 2024 年报数据：ROE、ROA、NIM、NPL、CAR
【输出要求】
  - figures/bank_roe_radar.pdf（5 行 × 5 维雷达图）
  - figures/bank_roe_heatmap.pdf（指标×银行 热力图）
  - 主题：theme_minimal(base_size=14)，Times 字体，colorblind-safe 配色
【环境检查】运行 /tmp/check_all_skills.sh
【原代码补充】R ggplot2 + fmsb 雷达图代码片段
【失败回退】如 econ-visualization 不可用，调用 paper-figure（Python matplotlib）
【预期结论】哪家银行综合表现最强？哪两家形成竞争梯队？
\end{lstlisting}

\textbf{学生练习}：将维度从 5 个扩展到 8 个（含拨备覆盖率/核心一级资本充足率/存款增速/贷款增速）。

\textbf{本节新增任务与前 5 个任务的对应关系}：
\begin{itemize}[leftmargin=1em]
  \item 任务 1（\ding{72}）：docx —— 文档办公类
  \item 任务 2（\ding{72}\ding{72}）：finance-expert —— 金融分析类
  \item 任务 3（\ding{72}\ding{72}\ding{72}）：xlsx —— 数据建模类
  \item 任务 4（\ding{72}\ding{72}\ding{72}）：自写 bank-risk-alert —— Skill 创建类
  \item 任务 5（\ding{72}\ding{72}\ding{72}）：Skill 安全审计 —— 合规审计类
  \item \textbf{任务 6（\ding{72}\ding{72}\ding{72}，新增）}：tushare-finance + china-stock-analysis —— 数据获取+金融分析类
  \item \textbf{任务 7（\ding{72}\ding{72}\ding{72}\ding{72}，新增）}：backtesting-trading-strategies —— 量化回测+金融分析类
  \item \textbf{任务 8（\ding{72}\ding{72}\ding{72}，新增）}：econ-visualization —— 学术图表+金融分析类
\end{itemize}

% ===========================================================
\section{学术与商业图形展示Skill（独立章节）}

% ===========================================================

\textbf{【本章特别要求】} 由于图形展示在银行金融场景中极度高频（监管报表、研报图表、学术论文图示），单独成节。涵盖5个核心Skill。

\subsection{econ-visualization（学术发表级图表，R ggplot2）}
\textbf{触发词}："出版级图表"、"发表级图"、"ggplot2图表"、"econ viz"。

\textbf{能力清单}：
\begin{itemize}[leftmargin=1em]
  \item 时间序列图、事件研究图、回归系数图、断点图、双轴图、热力图
  \item PDF/SVG矢量输出（300 DPI）
  \item 一致的学术主题（minimal, serif font, colorblind-safe）
\end{itemize}

\textbf{环境检查}：依赖R + ggplot2（见§R语言环境）。
\begin{lstlisting}[style=shell]
Rscript -e 'library(ggplot2); library(econvis); cat("OK")'
\end{lstlisting}

\textbf{提示词示例}：
\begin{lstlisting}
请使用 econ-visualization 技能生成以下学术图表：
- 数据：data/panel.csv（银行 2010-2024 年面板）
- 图表1：事件研究图（event_study），95% 置信区间阴影
- 图表2：四家银行 ROE 对比的双轴折线图
- 输出 PDF + PNG 至 figures/ch04_figures/
- 主题：theme_minimal，Times 字体，14pt
\end{lstlisting}

\subsection{paper-figure（论文图表生成器，Python matplotlib）}
\textbf{触发词}："画论文图"、"paper figure"、"matplotlib论文图表"。

\textbf{能力清单}：
\begin{itemize}[leftmargin=1em]
  \item 决策树自动选图（line/bar/scatter/heatmap/box/multi-panel）
  \item 共享样式脚本 \texttt{paper\_plot\_style.py}（DPI=300、PDF矢量、colorblind配色）
  \item 自动生成 \texttt{\textbackslash includegraphics} LaTeX片段
  \item 9类图表 × 8种样式矩阵
\end{itemize}

\textbf{环境检查}：依赖matplotlib + seaborn。
\begin{lstlisting}[style=shell]
pip install matplotlib seaborn
python3 -c "import matplotlib; print(matplotlib.__version__)"
\end{lstlisting}

\subsection{paper-illustration（论文插图/架构图，Gemini驱动）}
\textbf{触发词}："架构图"、"AI绘图"、"paper illustration"、"画架构图"。

\textbf{能力清单}：
\begin{itemize}[leftmargin=1em]
  \item 调用 \texttt{gemini-3-pro}优化布局 + \texttt{Paperbanana (gemini-3-pro-image-preview)}渲染
  \item CVPR/NeurIPS学术风格：圆角、渐变、4色配色、严格箭头方向
  \item 多轮Claude评分 ≥9分才输出
\end{itemize}

\textbf{环境检查}：
\begin{lstlisting}[style=shell]
export GEMINI_API_KEY='your-key'   # https://aistudio.google.com/app/apikey
ls qoder-skills-library/10-communication/paper-illustration/SKILL.md
\end{lstlisting}

\textbf{适用场景}：方法论流程图、模型架构图、概念图（\textbf{非统计数据图表}）。

\subsection{mermaid-diagram（Mermaid图表，文本驱动）}
\textbf{触发词}："流程图"、"时序图"、"甘特图"、"ER图"、"mermaid"。

\textbf{能力清单}：流程图、时序图、类图、状态图、甘特图、饼图、Gitgraph。

\textbf{环境检查}：
\begin{lstlisting}[style=shell]
npm install -g @mermaid-js/mermaid-cli
mmdc --version
\end{lstlisting}

\textbf{提示词示例}：
\begin{lstlisting}
请使用 mermaid-diagram 技能生成银行信贷审批流程图：
客户申请→征信查询→风控评分→人工审批→合同签订→放款
输出：figures/loan_workflow.mmd + 渲染 PNG
\end{lstlisting}

\subsection{学术图表与商业图表的差异表}

\begin{table}[H]
\centering
\begin{tabular}{lp{5cm}p{6cm}}
\toprule
\textbf{维度} & \textbf{学术图表（econ-visualization/paper-figure）} & \textbf{商业图表（mcp-server-chart）} \\
\midrule
目标期刊/受众 & 期刊审稿人（黑白打印友好） & 银行内部决策者（彩色屏幕优先） \\
字体 & Times/Computer Modern衬线体 & 微软雅黑/Arial无衬线体 \\
配色 & 2-4色，colorblind-safe & 品牌色，可渐变 \\
输出 & PDF/SVG矢量 & PNG/PPTX嵌入 \\
标题 & 写于LaTeX caption，图中无标题 & 图表内含大标题 \\
注释密度 & 极简，必要时footnote & 丰富，数据标签直接显示 \\
\bottomrule
\end{tabular}
\end{table}

% ===========================================================
\section{演示文稿Skill章节（独立章节）}

% ===========================================================

\textbf{【本章特别要求】} 演示文稿分学术与商业两条主线，对应不同Skill组合。涵盖6个核心Skill。

\subsection{学术演示文稿（Beamer LaTeX路线）}

\textbf{beamer-presentation Skill}：

\textit{触发词}："学术PPT"、"Beamer演示"、"学术汇报slides"。

\textit{能力清单}：metropolis/Madrid/默认主题；15/45/90分钟三档结构；演讲者备注。

\textit{输出}：\texttt{main.tex} + \texttt{main.pdf}。

\textit{适用}：开题答辩、学术会议汇报、job market talk。

\textbf{paper-slides Skill}：

\textit{触发词}："做会议幻灯片"、"conference slides"、"做slides"。

\textit{能力清单}（来自SKILL.md）：
\begin{itemize}[leftmargin=1em]
  \item 4种报告类型：poster-talk（3-5min, 5-8页）/ spotlight（5-8min, 8-12页）/ oral（15-20min, 15-22页）/ invited（30-45min, 25-40页）
  \item 7种会议配色方案（NeurIPS/ICML/ICLR/CVPR/ECCV/AAAI/ACL/EMNLP/GENERIC）
  \item 自动生成 \texttt{main.pdf}（Beamer）+ \texttt{presentation.pptx}（python-pptx）+ \texttt{TALK\_SCRIPT.md}（完整讲稿）
  \item 演讲者备注 + 预期Q\&A 8个
\end{itemize}

\textit{提示词示例}：
\begin{lstlisting}
/paper-slides "paper/" --talk_type: oral --venue: ICML --minutes: 20 --aspect: 16:9
\end{lstlisting}

\textbf{paper-poster Skill}：

\textit{触发词}："学术海报"、"poster"、"A0海报"。

\textit{能力清单}：A0/A1海报；多栏布局；QR码；图表嵌入。

\textit{适用}：学术会议海报展示。

\subsection{商业/银行演示文稿（PPTX路线）}

\textbf{pptx Skill}（通用PPT操作）：

\textit{触发词}："PPT"、"幻灯片"、"演示文稿"、"pitch deck"。

\textit{能力清单}（基于python-pptx + markitdown + LibreOffice）：
\begin{itemize}[leftmargin=1em]
  \item 文本提取（markitdown）、缩略图、XML解析
  \item 从零生成（pptxgenjs）或修改现有模板
  \item 9种专业配色（Deep Night、Woodland、Vibrant Coral、Earth Tone、Deep Sea、Graphite、Aquatic Trust、Wine \& Linen、Eucalyptus、Crimson Impact）
  \item 视觉设计原则：避免单调布局、避免默认蓝、加0.5" 留白、字号对比 ≥20pt
\end{itemize}

\textit{提示词示例}：
\begin{lstlisting}
请使用 pptx 技能生成《XX银行 2026 年战略汇报》PPT（20页）：
- 配色：Deep Night（navy + ice blue）
- 字体：Georgia（标题） + Calibri（正文）
- 每页含可视化元素（图标/图表/全屏图）
- 输出：reports/bank_strategy_2026.pptx
\end{lstlisting}

\textbf{qoderwork-ppt Skill}（智能模板生成）：

\textit{触发词}："智能PPT"、"QoderWork PPT"、"14模板"。

\textit{能力清单}：
\begin{itemize}[leftmargin=1em]
  \item 14种HTML模板自动匹配（含封面/分隔/内容/图标/全屏图等）
  \item 一键管线：\texttt{run-pipeline.js} → Puppeteer + dom-to-pptx转换
  \item 保留背景图、圆角、字体、全部样式
  \item Lucide图标库（本地静态资源，无需网络）
\end{itemize}

\textit{适用场景}：商业汇报、产品发布、培训材料（\textbf{需联网生成真实图片}）。

\textit{环境依赖}：Node.js + Puppeteer + LibreOffice + lucide-static（见§PPT演示文稿相关环境）。

\textbf{ppt-master Skill}（AI驱动多格式SVG内容生成）—— \textbf{【用户特别要求】}：

\textit{触发词}："ppt-master"、"SVG演示"、"AI幻灯片"。

\textit{能力清单}：
\begin{itemize}[leftmargin=1em]
  \item 多Agent协作：内容解析 → 视觉设计 → SVG渲染
  \item 输入支持：PDF / DOCX / URL / Markdown
  \item 输出：SVG页面 + 通过多角色协作导出PPTX
  \item 适用：把现成文档（论文/讲义/调研报告）一键转为专业幻灯片
\end{itemize}

\textit{环境检查}：
\begin{lstlisting}[style=shell]
export GEMINI_API_KEY='your_key'    # https://aistudio.google.com/app/apikey
export ANTHROPIC_API_KEY='sk-ant-...'
node --version && which libreoffice && which pdftoppm
\end{lstlisting}

\textit{与上述Skill的差异}：
\begin{itemize}[leftmargin=1em]
  \item \texttt{paper-slides}：论文 → 学术LaTeX slides
  \item \texttt{pptx}：通用PPT自由创作
  \item \texttt{qoderwork-ppt}：HTML模板 + 14套商业模板
  \item \textbf{ppt-master：多格式输入 + SVG中间层 + PPTX输出}（更智能的"文档→幻灯片"转换器）
\end{itemize}

\subsection{学术 vs 商业演示文稿Skill对照表}

\begin{table}[H]
\centering
\begin{tabular}{lll}
\toprule
\textbf{维度} & \textbf{学术（Beamer）} & \textbf{商业/银行（PPTX）} \\
\midrule
核心Skill & \texttt{beamer-presentation} / \texttt{paper-slides} & \texttt{pptx} / \texttt{qoderwork-ppt} / \texttt{ppt-master} \\
排版语言 & LaTeX（公式、参考文献无缝） & HTML→PPTX 或 PPTX直生成 \\
公式支持 & 完美（直接\$...\textbackslash\$） & 差（需图片或OMML） \\
品牌定制 & 弱（依赖Beamer主题） & 强（自由配色/字体） \\
可编辑性 & 弱（需懂LaTeX） & 强（PowerPoint直接改） \\
受众 & 学术评委、研讨会同行 & 行领导、客户、董事会 \\
典型产出 & \texttt{main.pdf} & \texttt{presentation.pptx} \\
\bottomrule
\end{tabular}
\end{table}

\subsection{提示词示例（综合）}

\begin{lstlisting}
请帮我制作一份《XX城商行 2026 年度数字化转型汇报》（30页 PPT）：
1. 使用 pptx 技能（商业风格，Deep Night 配色）
2. 同时使用 paper-slides 技能生成学术附录（Beamer 版本供发表）
3. 使用 paper-figure 技能生成关键图表（ROE 对比、贷款结构、客户增长）
4. 使用 mermaid-diagram 技能绘制数字化转型架构图
输出：
- 主报告：reports/digital_transformation_2026.pptx
- 学术附录：reports/digital_transformation_appendix.pdf
- 架构图：figures/digital_arch.png
\end{lstlisting}

% ===========================================================
\section{Skill编写方法论}
% ===========================================================

调用现有Skill只是起点，真正的能力提升来自于\textbf{编写自己的Skill}。本节系统讲解Skill编写的方法论，并通过实战案例带你完成第一个自编写Skill。

\subsection{Skill创建的元 Skill 工具箱}

在手工编写 SKILL.md 之前，强烈推荐先使用 \textbf{元 Skill} 工具箱。四个核心元 Skill 覆盖"起草→润色→发布→审计"全生命周期：

\begin{table}[H]
\centering
\begin{tabular}{lp{4cm}p{6cm}}
\toprule
\textbf{元 Skill} & \textbf{阶段} & \textbf{作用} \\
\midrule
\texttt{create-skill} & 起草 & 跟随交互向导生成完整 SKILL.md（含 frontmatter、argument-hint、allowed-tools、tags） \\
\texttt{humanizer-zh} & 润色 & 将 AI 生成腔的文本润色为自然语言，提升 Skill 可读性 \\
\texttt{write-skill-readme} & 发布 & 根据 SKILL.md 自动生成 README.md（含徽章、使用场景、示例代码） \\
\texttt{skill-auditor} & 审计 & 五维自动评分（触发词/输出/数据/合规/改进），输出结构化报告 \\
\bottomrule
\end{tabular}
\end{table}

\begin{lstlisting}
【Skill 创建 完整工作流提示词】

我要为本行创建一个反洗钱可疑交易报告生成 Skill（aml-suspicious-report），
请按以下流程协助：

Step 1（创建）：调用 create-skill 元 Skill 生成 SKILL.md 初始草稿
Step 2（参考）：调用 find-skills 查找同类 Skill（如 finance-expert、bank-risk-alert）作为参考
Step 3（润色）：调用 humanizer-zh 对草稿进行中文自然化润色，去除 AI 生成腔
Step 4（审计）：调用 skill-auditor 对草稿进行五维自动审计
Step 5（修复）：按审计反馈迭代修改（最多 2 轮）
Step 6（发布）：调用 write-skill-readme 生成 README.md
Step 7（诊断）：调用 install-skill-dependency 检查 Skill 是否能跑起来

【预期产出】
- aml-suspicious-report/SKILL.md（草稿定稿）
- aml-suspicious-report/README.md（自动生成）
- aml-suspicious-report/audit\_report.md（五维审计报告）
\end{lstlisting}

\begin{notebox}[元 Skill 工具箱的环境检查]
\begin{lstlisting}[style=shell]
# 检查 4 个元 Skill 都可访问
for skill in create-skill humanizer-zh write-skill-readme skill-auditor; do
  ls qoder-skills-library/19-devtools/$skill/SKILL.md 2>/dev/null \
    && echo "OK: $skill" \
    || echo "MISSING: $skill"
done
\end{lstlisting}
\end{notebox}

\subsection{触发词设计原则}

触发词是Skill与用户之间的"接口协议"，其设计质量直接决定Skill能否被正确调用。

\subsubsection{精确性vs覆盖性的平衡}

触发词设计面临一个核心矛盾：\textbf{太精确则覆盖不足，太宽泛则误触发}。

\begin{table}[H]
\centering
\begin{tabular}{lp{4cm}p{4cm}p{3cm}}
\toprule
\textbf{策略} & \textbf{优点} & \textbf{缺点} & \textbf{适用场景} \\
\midrule
精确触发词 & 几乎不误触发 & 可能遗漏部分调用场景 & 高风险操作（如交易执行） \\
宽泛触发词 & 覆盖面广 & 容易在无关场景被触发 & 低风险通用工具 \\
分层触发词 & 兼顾精确与覆盖 & 设计复杂度较高 & 专业领域Skill \\
\bottomrule
\end{tabular}
\end{table}

推荐采用\textbf{分层触发词}策略：设置核心触发词（高精确性）+ 扩展触发词（高覆盖性）+ 排除条件（防误触发）。

\subsubsection{中英文双语触发词设计}

在金融领域，许多术语存在中英文双语使用习惯。好的Skill应当同时覆盖两种语言：

\begin{lstlisting}
## When to Use

Use this skill when the user mentions:
- 中文触发词："信用风险评估""贷前审批""企业征信""风控分析"
- 英文触发词："credit risk", "credit assessment", "loan underwriting"
- 隐式触发：当用户提供了企业财务数据并要求"评估""打分""分析风险"时

Do NOT use this skill when:
- 用户只是查询某只股票的价格（应使用 stock-analysis）
- 用户要求进行投资组合优化（应使用 financial-modeling）
\end{lstlisting}

\subsubsection{误触发防护机制}

在\texttt{When to Use}中明确\textbf{排除条件}是防止误触发的有效手段。常见的排除场景包括：
\begin{itemize}[leftmargin=2em]
  \item 功能边界重叠：如"信用评估"vs"投资分析"都涉及财务数据处理
  \item 术语歧义：如"风险"可能是信用风险、市场风险或操作风险
  \item 用户意图不明确：如"帮我分析一下"——分析什么？
\end{itemize}

\begin{warningbox}[误触发的真实风险]
在银行业务场景中，Skill误触发可能导致严重后果。例如：
\begin{itemize}[leftmargin=1em]
  \item 本应调用"信用风险评估"Skill，却误调用了"投资分析"Skill，导致审批判断偏离合规标准
  \item 客户咨询理财建议时，误触发了"交易执行"Skill，导致非预期的交易指令
\end{itemize}
因此，金融类Skill的触发词设计应当偏保守，宁可漏触发（由用户明确指定），也不可误触发。
\end{warningbox}

\subsubsection{排除条件 vs Include 模式——两种触发词策略对照}

在金融领域，以下两种触发词设计策略各有适用场景，请根据业务特点选择：

\begin{table}[H]
\centering
\begin{tabular}{p{3.5cm}p{5cm}p{5cm}}
\toprule
\textbf{维度} & \textbf{排除条件模式（exclude-first）} & \textbf{Include 模式（include-first）} \\
\midrule
设计思路 & 先定义能力范围，再列出不适用场景 & 先列出适用场景，能力范围依场景推断 \\
适用场景 & 高风险、低错误容忍度（如交易执行、信用审批） & 低风险、探索性场景（如数据可视化、报告起草） \\
误触发风险 & 低（默认不触发，除非明显匹配） & 中（边缘场景可能误触发） \\
漏触发风险 & 中（用户需明确指定 Skill） & 低（场景描述完整即可触发） \\
金融场景示例 & \texttt{bank-risk-alert}：除非明确出现"风险预警"才触发 & \texttt{finance-expert}：只要涉及企业财务分析就可能触发 \\
\bottomrule
\end{tabular}
\end{table}

\textbf{案例对比}：同一业务场景"企业财务异常预警"，两种写法：

\begin{lstlisting}[style=json]
// === 排除条件模式（推荐金融风险场景）===
{
  "when_to_use": "处理企业财务异常预警任务",
  "include_keywords": ["财务异常", "财报预警", "财务预警"],
  "exclude_keywords": ["股价预警", "舆情预警", "市场风险"],
  "trigger_confidence_threshold": 0.85
}

// === Include 模式（适合探索性场景）===
{
  "when_to_use": "分析企业财务数据并发现异常",
  "include_keywords": [
    "财务分析", "财报解读", "异常检测", "三表分析",
    "财务预警", "杜邦分析", "现金流分析", "盈利质量"
  ],
  "exclude_keywords": [],
  "trigger_confidence_threshold": 0.6
}
\end{lstlisting}

\textbf{银行场景推荐}：金融业务优先采用\textbf{排除条件模式}，特别是涉及客户敏感数据、交易指令、风险评级等场景。

\subsection{输入输出规范}

\subsubsection{参数设计：必选vs可选}

Skill的输入参数应区分\textbf{必选参数}和\textbf{可选参数}：

\begin{lstlisting}
## 输入参数

### 必选参数（缺少则无法执行）
1. 客户名称 - 企业全称
2. 预警信号类型 - 枚举值：逾期/担保物贬值/关联企业异常/
   行业下行/资金异常流出
3. 风险等级 - 枚举值：黄色预警/橙色预警/红色预警

### 可选参数（缺失时使用默认值）
4. 贷款余额 - 默认：未提供
5. 到期日 - 默认：未提供
6. 客户经理 - 默认：当前登录用户
\end{lstlisting}

设计原则：\textbf{必选参数不超过4个}，否则会增加AI提取参数的失败率；可选参数提供合理默认值，降低用户的使用门槛。

\subsubsection{输出模板的结构化设计}

输出模板是Skill质量的直接体现。好的输出模板应当：
\begin{enumerate}[leftmargin=2em]
  \item \textbf{结构化}：使用标题、列表、表格等格式，而非大段文字
  \item \textbf{完整性}：覆盖该场景下用户需要的所有信息
  \item \textbf{可操作性}：输出不仅是分析结论，还应包含具体的行动建议
  \item \textbf{可审计}：关键结论有依据，便于复核
\end{enumerate}

\subsubsection{错误处理与降级策略}

Skill执行可能因各种原因失败（缺少依赖、输入不合法、外部服务不可用等）。好的Skill应定义明确的错误处理策略：

\begin{lstlisting}
## 错误处理

1. 若缺少必选参数，提示用户补充，而非猜测默认值
2. 若外部数据源不可用，使用本地缓存数据并标注"数据截止日期"
3. 若计算结果异常（如资产负债率 > 100%），输出警告而非静默忽略
4. 若无法完成全部步骤，已完成部分仍应输出，并标注"未完成步骤"
\end{lstlisting}

\subsection{实战：编写bank-risk-alert Skill（\ding{72}\ding{72}\ding{72}）}

\textbf{场景}：你是银行科技部的AI工程师，需要为全行客户经理打造一个「贷后风险预警话术生成器」Skill。

\begin{enumerate}[leftmargin=2em]
  \item 新建文件 \texttt{skills/bank-risk-alert/SKILL.md}，要求包含：
\end{enumerate}

\begin{lstlisting}
请帮我编写一个 bank-risk-alert Skill（SKILL.md），要求如下：

【Skill 基本信息】
- name: bank-risk-alert
- description: 银行贷后风险预警话术生成器，根据客户风险信号自动生成
  分级预警通知和客户经理跟进话术
- 触发词："风险预警""贷后预警""逾期提醒""risk alert"

【输入参数】
1. 客户名称
2. 预警信号类型（逾期/担保物贬值/关联企业异常/行业下行/资金异常流出）
3. 风险等级（黄色预警/橙色预警/红色预警）
4. 贷款余额与到期日

【输出模板（必须包含）】
1. 预警摘要卡片（一屏可看完的关键信息）
2. 客户经理电话话术（开场白→风险告知→解决方案→行动承诺→结束语）
3. 内部上报邮件模板（收件人=分管行长，含风险描述+处置建议+时间表）
4. 跟进工单（任务项+负责人+截止日期+检查标准）

【分级逻辑】
- 黄色：关注但不急，7日内跟进
- 橙色：需立即电话联系，3日内上报
- 红色：启动应急预案，当日上报 + 法务介入

【示例输入输出】
提供至少 2 个完整的输入-输出示例

【Best Practices】
至少 3 条银行贷后管理最佳实践

请确保 SKILL.md 格式规范，能被 AI IDE 正确识别和调用。
\end{lstlisting}

\textbf{完成后测试}：在TRAE/Qoder中加载该Skill，输入测试提示词验证：

\begin{lstlisting}
请生成风险预警：客户"恒达机械"出现连续2期利息逾期，
贷款余额500万，到期日2026-09-30，预警等级为橙色。
\end{lstlisting}

\textbf{拓展C：自写bank-product-recommender Skill}

高阶学生可额外编写一个「银行产品智能推荐」Skill：

\begin{itemize}[leftmargin=2em]
  \item 输入：客户画像（年龄/收入/风险偏好/持有产品）
  \item 输出：Top 3推荐产品 + 推荐理由 + 客户经理话术
  \item 逻辑：保守型→大额存单/国债；稳健型→固收+理财；进取型→基金/结构性存款
\end{itemize}

\subsection{Skill测试与迭代方法}

编写完Skill只是第一步，确保其可靠运行需要系统化的测试与迭代。

\subsubsection{测试维度}

一个Skill的测试应覆盖以下维度：

\begin{table}[H]
\centering
\begin{tabular}{lp{8cm}}
\toprule
\textbf{测试维度} & \textbf{测试方法} \\
\midrule
触发准确性 & 输入相关/无关提示词，验证是否正确触发/不触发 \\
参数提取 & 输入包含不同参数组合的提示词，验证AI是否正确提取 \\
输出完整性 & 对比输出模板，检查每个必填项是否都被覆盖 \\
边界情况 & 输入极端值（如贷款余额为0、风险等级不合法），验证错误处理 \\
跨语言 & 分别用中文和英文提示词调用，验证双语支持 \\
\bottomrule
\end{tabular}
\end{table}

\subsubsection{迭代优化流程}

\begin{practicebox}[Skill迭代优化四步法]
\begin{enumerate}[leftmargin=1em]
  \item \textbf{记录问题}：每次使用中遇到的问题（误触发、输出遗漏、格式错误等）记录到Skill的Changelog中
  \item \textbf{定位原因}：分析问题出在哪个环节——触发词不够精准？Instructions步骤遗漏？Best Practices提示不足？
  \item \textbf{定向修复}：只修改问题相关的部分，避免"大改"引入新问题
  \item \textbf{回归测试}：修复后重新运行之前的测试用例，确保没有引入新问题
\end{enumerate}
\end{practicebox}

\begin{notebox}[版本管理建议]
建议为Skill维护一个Changelog，记录每次修改的内容和原因：
\begin{lstlisting}
## Changelog

### v1.1.0 (2026-06-08)
- 新增排除条件：区分"信用风险"与"市场风险"
- 修复：橙色预警的跟进时限从"5日"更正为"3日"
- 新增：抵押物贬值场景的输出模板

### v1.0.0 (2026-05-15)
- 初始版本发布
\end{lstlisting}
\end{notebox}

% ===========================================================
\section{Skill安全审计与合规评估}
% ===========================================================

银行在引入外部AI Skill前，必须进行严格的安全与合规审计。这不仅关乎数据安全，更是监管合规的硬性要求。

\subsection{自动化审计工具：skill-auditor 元 Skill}

\textbf{【本节特别推荐】} 传统的人工五维审计耗时较长，本项目提供 \texttt{skill-auditor} 元 Skill，可对任意 SKILL.md 进行自动化五维评分。

\subsubsection{skill-auditor 的能力清单}

\begin{itemize}[leftmargin=1em]
  \item 自动检查 触发词准确性（中英文覆盖、排除条件、误触发风险）
  \item 自动检查 输出完整性（Instructions 步骤、输出模板、示例输入输出）
  \item 自动检查 数据依赖（明文敏感词、必填字段、数据脱敏规则）
  \item 自动检查 合规性（免责声明、监管提示、审批留痕）
  \item 自动生成 改进建议（按优先级排序、可直接执行）
\end{itemize}

\subsubsection{3 步环境检查}

\begin{lstlisting}[style=shell]
# Step 1：路径检查（skill-auditor 是 Skill 文件，无独立二进制）
ls qoder-skills-library/19-devtools/skill-auditor/SKILL.md
# Step 2：Skill 可访问性
qoder skill ls | grep skill-auditor
# Step 3：调用测试
qoder skill run skill-auditor --target=.agents/skills/bank-risk-alert/SKILL.md
\end{lstlisting}

\subsubsection{调用提示词}

\begin{lstlisting}
请调用 skill-auditor 技能对 .agents/skills/bank-risk-alert/SKILL.md
进行五维自动审计：

【调用 Skill】skill-auditor、security-best-practices
【输入】.agents/skills/bank-risk-alert/SKILL.md
【输出】reports/audit_bank-risk-alert.md（含五维评分表 + 改进建议）
【审计依据】security-best-practices Skill 提供的银行 Skill 安全清单
\end{lstlisting}

\subsubsection{人工审计 vs skill-auditor 自动审计 对比}

\begin{table}[H]
\centering
\begin{tabular}{p{3.5cm}p{5cm}p{5cm}}
\toprule
\textbf{维度} & \textbf{人工审计} & \textbf{skill-auditor 自动审计} \\
\midrule
耗时 & 2-4 小时/Skill & 30 秒/Skill \\
一致性 & 受审计员经验影响大 & 标准化输出，跨 Skill 可比 \\
覆盖度 & 受审计员精力限制可能漏项 & 覆盖所有 checklist 项目 \\
深度 & 能识别语义上下文问题 & 只能识别模式化问题（关键词、结构） \\
主观判断 & 能结合业务场景灵活判断 & 仅按预设规则评分 \\
监管证据 & 需审计员签字 & 自动生成审计日志，可追溯 \\
适用场景 & 重要 Skill 上线前终审 & 全量 Skill 定期巡检 \\
\bottomrule
\end{tabular}
\end{table}

\textbf{推荐组合}：先用 skill-auditor 完成全量 Skill 的初筛（<1 小时），再对高风险 Skill 进行人工复审。

\subsection{外部 Skill 白名单/黑名单管理}

银行在引入外部 Skill 时，应建立\textbf{分级管理}制度：

\begin{table}[H]
\centering
\begin{tabular}{lp{3cm}p{5cm}p{4cm}}
\toprule
\textbf{类别} & \textbf{准入要求} & \textbf{典型 Skill} & \textbf{数据流向} \\
\midrule
白名单（开放） & 官方认证 + 内部审计通过 & \texttt{docx}、\texttt{xlsx}、\texttt{mermaid-diagram} & 仅本地处理 \\
灰名单（审批） & 业务部门申请 + 合规复审 & \texttt{finance-expert}、\texttt{tushare-finance} & 可联网但禁止传敏感数据 \\
黑名单（禁用） & 任何情况下禁止使用 & 未审计的第三方 Skill & 禁止传输任何数据 \\
\bottomrule
\end{tabular}
\end{table}

\begin{lstlisting}[style=json]
// 推荐的白名单管理文件：.agents/skill-whitelist.json
{
  "whitelist": [
    {"name": "docx", "version": ">=1.0.0", "approved_by": "科技部", "approved_at": "2025-12-01"},
    {"name": "xlsx", "version": ">=1.0.0", "approved_by": "科技部", "approved_at": "2025-12-01"},
    {"name": "econ-visualization", "version": ">=1.2.0", "approved_by": "风控部", "approved_at": "2026-01-15"}
  ],
  "graylist": [
    {"name": "tushare-finance", "constraint": "禁止传入客户身份证/账号", "approved_by": "数据治理委员会"}
  ],
  "blacklist": [
    {"name": "未审计的任意 npm skills", "reason": "供应链风险"}
  ]
}
\end{lstlisting}

\subsection{银行业AI应用合规框架}

\subsubsection{核心法规依据}

当前中国银行业AI应用需要遵循的法规框架包括：

\begin{enumerate}[leftmargin=2em]
  \item \textbf{《银行保险机构信息科技外包风险管理办法》}（银保监办发〔2021〕141号）：将AI工具纳入信息科技外包管理范畴，要求对外包服务商进行风险评估、合同约束和持续监控。使用外部Skill本质上是一种"知识外包"，需要纳入此框架管理。
  \item \textbf{《生成式人工智能服务管理暂行办法》}（2023年8月施行）：要求提供生成式AI服务的组织进行安全评估、算法备案，并对生成内容标注标识。银行作为AI服务的使用方，需确保所用AI工具的提供方已完成合规手续。
  \item \textbf{《商业银行数据安全管理办法》}（征求意见稿）：要求数据处理活动进行安全影响评估，敏感数据出境需经审批。AI Skill在处理客户数据时，必须符合数据分类分级保护要求。
\end{enumerate}

\subsubsection{金融数据分类分级}

根据监管要求，金融数据按敏感程度分为三个级别：

\begin{table}[H]
\centering
\begin{tabular}{lp{4cm}p{6cm}}
\toprule
\textbf{级别} & \textbf{数据类型} & \textbf{Skill处理要求} \\
\midrule
一般数据 & 公开的银行产品信息、行业统计 & 可正常处理，无特殊限制 \\
重要数据 & 客户姓名、资产规模、交易频次 & 必须脱敏处理，禁止明文传输至外部AI \\
核心数据 & 身份证号、账户密码、交易明细 & 禁止输入AI系统，仅允许本地部署模型处理 \\
\bottomrule
\end{tabular}
\end{table}

\begin{warningbox}[数据脱敏实操要点]
在使用Skill处理客户数据时，务必遵循以下脱敏规则：
\begin{itemize}[leftmargin=1em]
  \item 客户姓名：用"张**"替代全名
  \item 身份证号：仅保留前3位和后4位，如"510****1234"
  \item 账号信息：仅保留后4位，如"****5678"
  \item 金额数据：可保留具体数值用于计算，但禁止与客户身份信息关联输出
  \item 联系方式：禁止输入AI系统，在输出模板中用"[客户电话]"占位符替代
\end{itemize}
\end{warningbox}

\subsubsection{AI输出的法律责任边界}

AI Skill生成的分析报告、风险评估结论等，在法律上如何定性？这是一个尚未完全明确但必须关注的问题：

\begin{itemize}[leftmargin=2em]
  \item \textbf{辅助决策vs自动决策}：如果AI输出仅作为参考，最终决策由人工做出，则责任主体为人；如果AI输出直接触发操作（如自动拒绝贷款），则可能构成自动决策，责任归属更为复杂。
  \item \textbf{投资建议vs分析报告}：如果Skill输出被解读为投资建议，则可能违反《证券法》关于投资咨询资质的要求；如果是分析报告并附有免责声明，则合规风险较低。
  \item \textbf{可解释性义务}：根据《个人信息保护法》第24条，通过自动化决策方式作出对个人权益有重大影响的决定时，个人信息处理者应同时提供不予接受的选项，并说明决策机制。
\end{itemize}

\subsection{五维审计方法论}

\subsubsection{审计维度详解}

对金融Skill的安全审计应从以下五个维度系统展开：

\begin{table}[H]
\centering
\begin{tabular}{clp{6.5cm}}
\toprule
\textbf{\#} & \textbf{审计维度} & \textbf{检查要点} \\
\midrule
1 & 触发词设计 & 是否精准？是否会误触发？中英文覆盖？ \\
2 & 输出质量 & 步骤是否完整？是否有遗漏场景？ \\
3 & 数据安全 & 是否可能泄露客户隐私？是否有数据脱敏机制？ \\
4 & 合规风险 & 输出是否可能构成投资建议？是否有免责声明？ \\
5 & 改进建议 & 至少提出3条具体的改进方案 \\
\bottomrule
\end{tabular}
\end{table}

各维度的详细检查清单：

\begin{enumerate}[leftmargin=2em]
  \item \textbf{触发词设计}（10分制）：
    \begin{itemize}
      \item 核心触发词是否与Skill功能精确对应？
      \item 是否覆盖了用户可能使用的同义表述？
      \item 是否有明确的排除条件防止误触发？
      \item 中英文触发词是否齐全？
    \end{itemize}
  \item \textbf{输出质量}（10分制）：
    \begin{itemize}
      \item Instructions步骤是否完整、有序？
      \item 是否覆盖了该领域的主要场景？
      \item 输出模板是否结构化、可操作？
      \item 是否有示例输入输出供参考？
    \end{itemize}
  \item \textbf{数据安全}（10分制）：
    \begin{itemize}
      \item 是否明确标注了禁止输入的数据类型？
      \item 输出中是否会包含敏感信息的明文？
      \item 是否有数据脱敏的具体规则？
      \item Skill执行是否涉及数据传输至外部服务？
    \end{itemize}
  \item \textbf{合规风险}（10分制）：
    \begin{itemize}
      \item 输出是否附有免责声明？
      \item 是否可能被解读为投资建议或审批决定？
      \item 是否符合银行数据分级保护要求？
      \item 是否有审计日志或操作记录机制？
    \end{itemize}
  \item \textbf{改进建议}（10分制）：
    \begin{itemize}
      \item 改进方案是否具体可操作？
      \item 是否考虑了实施成本与收益？
      \item 是否按优先级排序？
    \end{itemize}
\end{enumerate}

\subsubsection{审计报告提示词}

\begin{lstlisting}
请对 finance-expert Skill 进行银行级安全审计：

1. 阅读该 Skill 的完整 SKILL.md 内容
2. 从"触发词精准性/输出完整性/数据安全/合规风险/可改进空间"五个维度评估
3. 特别关注：
   - 该 Skill 是否会在不合适的场景被误触发？
   - 输出结果是否可能被视为正式投资建议（违反银保监规定）？
   - 是否有机制防止客户真实数据泄露到 AI 模型？
4. 输出结构化审计报告（>=300字），含评分（每维度10分制）+ 改进建议

保存到 reports/skill_audit_report.md
\end{lstlisting}

\subsubsection{审计报告模板}

审计报告应遵循以下标准模板：

\begin{lstlisting}
# 金融Skill安全审计报告

## 基本信息
- 审计对象：[Skill名称] v[版本号]
- 审计日期：[YYYY-MM-DD]
- 审计人员：[姓名/工号]
- 审计依据：《银行保险机构信息科技外包风险管理办法》等

## 维度一：触发词设计 [X/10]
### 发现
[具体描述触发词设计的优点和问题]
### 风险等级：[高/中/低]

## 维度二：输出质量 [X/10]
### 发现
[具体描述输出质量的评估结果]
### 风险等级：[高/中/低]

## 维度三：数据安全 [X/10]
### 发现
[具体描述数据安全方面的评估结果]
### 风险等级：[高/中/低]

## 维度四：合规风险 [X/10]
### 发现
[具体描述合规风险的评估结果]
### 风险等级：[高/中/低]

## 维度五：改进建议 [X/10]
### 改进方案
1. [具体方案1] - 优先级：[高/中/低]
2. [具体方案2] - 优先级：[高/中/低]
3. [具体方案3] - 优先级：[高/中/低]

## 综合评定
- 总分：[X/50]
- 审计结论：[通过/有条件通过/不通过]
- 建议措施：[具体建议]
\end{lstlisting}

\begin{casebox}[金融Skill的银保监合规审查案例]
\textbf{案例背景}：某城商行科技部引入了一个名为\texttt{smart-loan-advisor}的外部Skill，用于辅助客户经理进行贷款产品推荐。在上线前，合规部门对其进行了审查。

\textbf{审查发现}：

\textbf{问题1——触发词过于宽泛}：该Skill的触发词包含"贷款""利率""还款"等通用词汇，导致客户在咨询存款利率时也触发了贷款推荐Skill，构成不当营销。

\textbf{问题2——输出缺乏免责声明}：Skill生成的推荐报告未标注"本报告仅供参考，不构成贷款承诺"，可能被客户解读为贷款审批结果。

\textbf{问题3——数据未脱敏}：Skill的Example Prompts中使用了真实的客户姓名和身份证号作为示例，违反了《个人信息保护法》。

\textbf{处理结果}：合规部门要求该Skill在完成以下整改后方可上线——（1）缩小触发词范围至"贷款推荐""产品对比"等；（2）在输出模板末尾增加免责声明；（3）替换示例中的真实个人信息为脱敏数据。

\textbf{教训}：金融Skill的合规审查不是可选项，而是上线前的必经环节。一份完整的审计报告可以在银保监检查时作为合规履职的证据。
\end{casebox}

% ===========================================================
\section{进阶选做}
% ===========================================================

\subsection{金融新闻舆情分析（\ding{72}\ding{72}\ding{72}\ding{72}）}

\textbf{场景}：银行舆情监控部门需要每日跟踪银行业相关新闻，进行情绪分析和风险预警。

\textbf{技能安装}：

\begin{lstlisting}[style=shell]
npx skills add sundial-org/awesome-openclaw-skills@finance-news -y
\end{lstlisting}

\textbf{操作提示词}：

\begin{lstlisting}
请使用 finance-news 技能完成银行业舆情分析任务：

【第 1 步】新闻采集
搜索近一周中国银行业相关重大新闻，来源包括：
- 央行/银保监/金融监管总局政策发布
- 上市银行公告（如业绩快报、高管变动、处罚信息）
- 主流财经媒体报道（第一财经、证券时报、21世纪经济报道）

【第 2 步】情绪评分
对每条新闻进行结构化评分：
| 新闻标题 | 来源 | 日期 | 情绪(+1/0/-1) | 影响程度 | 涉及机构 |

【第 3 步】风险信号识别
从新闻中提取潜在风险信号：
- 监管处罚 → 合规风险
- 房地产相关 → 资产质量风险
- 利率政策变化 → 利差收窄风险
- 科技公司跨界 → 竞争风险

【第 4 步】生成舆情简报
输出不超过 500 字的银行业周度舆情简报，包含：
1. 本周要闻 TOP 5
2. 情绪倾向总结（乐观/中性/悲观）
3. 需关注的风险点
4. 对本行业务的潜在影响

保存到 reports/bank_sentiment_weekly.md
\end{lstlisting}

\subsection{咨询框架分析银行数字化转型（\ding{72}\ding{72}\ding{72}\ding{72}）}

\textbf{技能安装}：

\begin{lstlisting}[style=shell]
npx skills add aznatkoiny/zai-skills@consulting-frameworks -y
\end{lstlisting}

\textbf{操作提示词}：

\begin{lstlisting}
请使用 consulting-frameworks 技能，用波特五力模型分析
中国银行业在 AI 时代面临的竞争格局变化：

1. 现有竞争者（国有大行 vs 股份行 vs 城农商行）
2. 潜在进入者（互联网银行、金融科技公司）
3. 替代品威胁（支付宝/微信支付、数字人民币）
4. 供应商议价力（科技供应商、数据供应商）
5. 客户议价力（零售客户、对公客户）

要求：
- 每个力量至少分析 3 个要点
- 结合 2025-2026 年的实际政策与市场变化
- 最后给出：中小银行应如何利用 AI 建立差异化竞争优势？
- 输出 Markdown 报告 + 五力模型的 mermaid 图

保存到 reports/porter_five_forces_banking.md
\end{lstlisting}

\subsection{financial-modeling Skill：银行 DCF 估值（\ding{72}\ding{72}\ding{72}\ding{72}）}

\textbf{场景}：本行拟对一家拟授信客户进行估值，作为贷款额度的参考依据。

\textbf{技能安装}：

\begin{lstlisting}[style=shell]
ls qoder-skills-library/12-finance/financial-modeling/SKILL.md
# 验证 Python 环境
python3 -c "import numpy_financial; print('OK')" || pip install numpy-financial
\end{lstlisting}

\textbf{3 步环境检查}：
\begin{lstlisting}[style=shell]
which python3
python3 -c "import numpy_financial; print(numpy_financial.__version__)"
ls qoder-skills-library/12-finance/financial-modeling/SKILL.md
\end{lstlisting}

\textbf{操作提示词（Skill 化包装）}：

\begin{lstlisting}
请调用 financial-modeling 技能完成商业银行 DCF 估值：

【调用 Skill】financial-modeling、tushare-finance
【背景】拟授信客户 X 银行 DCF 估值
【输入数据】
  - 拉取 X 银行 2018-2024 年报数据：净利润、净资产、ROE、增长
  - 折现率 WACC 假设：股权成本 10%、债务成本 4.5%、资本结构 8% 股权 92% 债务
  - 增长率阶段：明示期 5 年 + 永续增长 2.5%
【输出要求】
  - reports/bank_dcf_valuation.xlsx（含 3 表：FCFE/FCFF/股息贴现）
  - figures/dcf_sensitivity.pdf（折现率与永续增长率敏感性分析）
  - reports/dcf_report.md（含估值结论与风险提示）
【环境检查】运行 /tmp/check_all_skills.sh
【原代码补充】DCF 标准代码片段（numpy_financial.npv / irr）
【失败回退】如 Skill 不可用，手写三表 DCF 逻辑
【预期结论】X 银行每股内在价值区间、对应市净率
\end{lstlisting}

\textbf{学生练习}：调整 WACC 和增长率假设，观察估值结果的敏感性。

\subsection{humanizer-zh Skill：AI 文本检测与改写（\ding{72}\ding{72}\ding{72}）}

\textbf{场景}：撰写信贷调查报告时，需要去除 AI 生成的“机械感”，提升文本自然度。

\textbf{技能安装}：

\begin{lstlisting}[style=shell]
ls qoder-skills-library/07-writing/humanizer-zh/SKILL.md
# humanizer-zh 一般仅需联网调用 AI 检测 API
curl -I https://api.openai.com/ | head -1
export OPENAI_API_KEY='sk-...'
\end{lstlisting}

\textbf{3 步环境检查}：
\begin{lstlisting}[style=shell]
ls qoder-skills-library/07-writing/humanizer-zh/SKILL.md
[ -n "$OPENAI_API_KEY" ] && echo "API_KEY SET" || echo "MISSING"
\end{lstlisting}

\textbf{操作提示词（Skill 化包装）}：

\begin{lstlisting}
请调用 humanizer-zh 技能对以下信贷调查报告进行人性化润色：

【调用 Skill】humanizer-zh
【输入】reports/credit_risk_hengda.md
【输出要求】
  - reports/credit_risk_hengda_humanized.md（润色后版本）
  - 检测报告：AI 痕迹评分（越低越自然）、修改点列表
【质量要求】不能修改专业术语与数据，口语化调整需保留“银行内部报告”严肃风格
\end{lstlisting}

\textbf{学生练习}：对比润色前后文本的 AI 检测评分（如使用 GPTZero）。

\subsection{consulting-frameworks Skill：银行业 SWOT + BCG 矩阵（\ding{72}\ding{72}\ding{72}）}

\textbf{场景}：本行战略部需要制定 2026-2030 年数字化战略，需用多种咨询框架进行多维度分析。

\textbf{技能安装}：

\begin{lstlisting}[style=shell]
ls qoder-skills-library/13-consulting/consulting-frameworks/SKILL.md
# 一般仅需 LLM 调用能力
\end{lstlisting}

\textbf{3 步环境检查}：
\begin{lstlisting}[style=shell]
ls qoder-skills-library/13-consulting/consulting-frameworks/SKILL.md
\end{lstlisting}

\textbf{操作提示词（Skill 化包装）}：

\begin{lstlisting}
请调用 consulting-frameworks 技能完成商业银行多框架战略分析：

【调用 Skill】consulting-frameworks
【背景】XX 城商行 2026-2030 年数字化战略制定
【输入】本行近 3 年财务、资产质量、客户结构数据
【输出要求】
  - reports/swot_matrix.md（4×4 矩阵 + 战略选择象限图）
  - reports/bcg_portfolio.md（业务条线 BCG 矩阵：明星/现金牛/问题/瘦狗）
  - reports/strategy_summary.md（多框架交集分析出的 TOP3 战略选择）
【图表要求】每份报告配 1-2 个 mermaid 架构图
【环境检查】运行 /tmp/check_all_skills.sh
【原代码补充】SWOT/BCG 框架核心问题列表
【失败回退】如 Skill 不可用，手动按框架填写
【预期结论】本行处于哪象限？数字化投资应集中在哪三个领域？
\end{lstlisting}

\textbf{学生练习}：选取本省某城商行作为案例，完成完整 SWOT + BCG 分析（≥1500 字报告）。

\textbf{进阶选做 5 个任务的对应 Skill 一览}：
\begin{itemize}[leftmargin=1em]
  \item 任务 1：\texttt{finance-news} + \texttt{news-read}（舆情分析）
  \item 任务 2：\texttt{consulting-frameworks}（波特五力）
  \item \textbf{任务 3（新增）}：\texttt{financial-modeling}（DCF 估值）
  \item \textbf{任务 4（新增）}：\texttt{humanizer-zh}（AI 文本检测与改写）
  \item \textbf{任务 5（新增）}：\texttt{consulting-frameworks}（SWOT + BCG 矩阵）
\end{itemize}

% ===========================================================
\section*{本章小结}
% ===========================================================

\addcontentsline{toc}{section}{本章小结}

本章围绕"如何赋予AI金融专业能力"这一核心问题，从设计原理、生态全景、调用实战、编写方法、安全审计五个层面系统讲解了Skill体系。

\begin{itemize}[leftmargin=2em]
  \item \textbf{设计原理}：Skill是知识工程从专家系统到AI Skill演进的产物，以Markdown文档承载领域知识，遵循"发现--理解--执行--保障"的逻辑架构。核心口诀是"MCP给能力，Skill给方法"。
  \item \textbf{生态全景}：当前金融Skill已覆盖数据获取、专业分析、文档生成、咨询框架四大类别，通过skills.sh市场或本地目录进行安装和管理。
  \item \textbf{调用实战}：通过docx（客户回访报告）、finance-expert（信用风险评估）、xlsx（贷款FTP定价模型）三个任务，由浅入深掌握了Skill的调用方法。
  \item \textbf{编写方法}：掌握了触发词设计的精确性与覆盖性平衡、输入输出规范、错误处理策略，并通过编写bank-risk-alert Skill进行了实战演练。
  \item \textbf{安全审计}：从银行业合规框架出发，建立了五维审计方法论（触发词/输出/数据/合规/改进），确保AI Skill在金融场景中的安全合规使用。
\end{itemize}

\keyterms{Skill、SKILL.md、触发词、MCP vs Skill、知识工程、FTP转移定价、五C模型、RFM分群、数据脱敏、合规审计、五维审计法、金融数据分类分级、Skill库全景、167个Skill、20分类、元Skill、create-skill、find-skills、skill-auditor、write-skill-readme、install-skill-dependency、luban、tushare-finance、china-stock-analysis、backtesting-trading-strategies、financial-modeling、econ-visualization、ppt-master、qoderwork-ppt、paper-slides、paper-figure、mermaid-diagram、white-list/black-list、exclude-first/include-first、环境准备速查表、R环境、Stata路径、Python环境、Node.js、LualOffice、poppler-utils}

\section*{思考题}
\addcontentsline{toc}{section}{思考题}

\begin{enumerate}[leftmargin=2em]
  \item 如果银行需要同时引入"信用风险评估"和"投资组合分析"两个Skill，它们的触发词应如何设计才能避免冲突？请给出具体的触发词方案和排除条件。

  \item 从知识工程的角度分析，Skill体系相比传统规则引擎（如Drools）和机器学习模型，在金融风控场景中各自的优势与局限是什么？

  \item 某银行计划将AI Skill用于自动化的贷后风险预警，要求橙色及以上预警自动发送通知给客户经理。从《银行保险机构信息科技外包风险管理办法》和《个人信息保护法》的角度，分析该方案需要满足哪些合规条件？

  \item 在FTP定价模型中，如果央行为应对经济下行而下调LPR，银行应如何调整定价参数？请分析FTP、风险溢价、目标利润三个参数的调整逻辑及其对最终贷款利率的影响。

  \item 设计一个"银行反洗钱预警"Skill的SKILL.md大纲，包括name、触发词、输入参数、输出模板、分级逻辑和Best Practices。特别说明：该Skill如何处理可疑交易报告（STR）的合规要求？

  \item \textbf{【新增】} 在本机 167 个 Skill / 20 分类体系中，如果要为"中小企业信用贷款审批"场景找到最合适的 Skill 组合，你会如何选择？请说明选择策略并列出 5-8 个 Skill 及其协同方式。

  \item \textbf{【新增】} 元 Skill 体系（create-skill / find-skills / skill-auditor / luban）与普通业务 Skill 在使用场景上有什么本质区别？请举例说明元 Skill 如何在 Skill 创建 → 发现 → 审计 → 打磨全生命周期中发挥作用。

  \item \textbf{【新增】} 在金融场景中，是应该使用排除条件模式（exclude-first）还是 Include 模式（include-first）？请从误触发风险、漏触发风险两个角度对比，并给出你认为应该采用哪种模式的判断准则。

  \item \textbf{【新增】} 假设你计划同时使用 tushare-finance、backtesting-trading-strategies、econ-visualization 三个 Skill 完成"银行股估值+回测+可视化"项目，请说明：(a) 三个 Skill 的调用顺序与依赖关系；(b) 每个 Skill 需要的环境检查步骤；(c) 如何通过 \texttt{/tmp/check_all_skills.sh} 一键验证环境。

  \item \textbf{【新增】} skill-auditor 自动化审计与人工五维审计各有优劣。请说明：什么场景下应该以 skill-auditor 初筛为主、人工复审为辅？什么场景下必须人工审计？
\end{enumerate}

\section*{实验报告要求}
\addcontentsline{toc}{section}{实验报告要求}

\begin{table}[H]
\centering
\begin{tabular}{cll}
\toprule
\textbf{\#} & \textbf{交付物} & \textbf{形式} \\
\midrule
1 & 客户回访报告 & reports/客户回访报告\_C00001.docx \\
2 & 信用风险分析报告 & reports/credit\_risk\_hengda.md \\
3 & 贷款定价模型 & reports/loan\_pricing\_model.xlsx \\
4 & 自写bank-risk-alert Skill & skills/bank-risk-alert/SKILL.md \\
5 & bank-risk-alert测试截图 & 输入 + 输出截图 \\
6 & 金融Skill审计报告 & reports/skill\_audit\_report.md \\
7 & 实验心得（$\geq$ 200字） & 含AI出错案例反思 \\
8 & \textbf{【新增】} 银行股估值对比表 & reports/bank\_valuation\_compare.xlsx \\
9 & \textbf{【新增】} 量化回测报告 & reports/strategy\_backtest\_results.xlsx \\
10 & \textbf{【新增】} ROE 雷达图 & figures/bank\_roe\_radar.pdf \\
\bottomrule
\end{tabular}
\end{table}

\textbf{加分项}（非必做，完成可额外加分）：

\begin{itemize}[leftmargin=2em]
  \item 完成「银行客户RFM分群」或「舆情分析」任务（+5分）
  \item 自写bank-product-recommender Skill并通过测试（+5分）
  \item 完成「波特五力银行数字化转型」分析（+5分）
  \item \textbf{【新增】} 完成 tushare-finance A/H 股估值对比（任务 6）（+5分）
  \item \textbf{【新增】} 完成 backtesting-trading-strategies 三策略回测（任务 7）（+5分）
  \item \textbf{【新增】} 完成 econ-visualization ROE 雷达图（任务 8）（+5分）
  \item \textbf{【新增】} 完成 financial-modeling DCF 估值（任务 3）（+5分）
  \item \textbf{【新增】} 完成 humanizer-zh AI 文本改写（任务 4）（+5分）
  \item \textbf{【新增】} 完成 consulting-frameworks SWOT+BCG 分析（任务 5）（+5分）
\end{itemize}
